\documentclass[11pt,a4paper,landscape,twocolumn]{article}
\usepackage{fontspec}
\usepackage{polyglossia}
\setdefaultlanguage{russian}\setotherlanguage{english}
% Шрифты лежат рядом, в папке fonts/, и подключаются по имени файла: устанавливать их в систему не нужно.
% Чтобы взять системный шрифт, замените строку целиком, например: \setmainfont{Times New Roman}\newfontfamily\cyrillicfont{Times New Roman}
\newcommand{\optserif}{Path=fonts/,Extension=.ttf,UprightFont=*-Regular,BoldFont=*-Bold,ItalicFont=*-Italic,BoldItalicFont=*-BoldItalic}
\newcommand{\optsans}{Path=fonts/,Extension=.otf,UprightFont=*-Regular,BoldFont=*-ExtraBold,ItalicFont=*-Italic,BoldItalicFont=*-ExtraBoldItalic,Scale=0.92}
\newcommand{\optmono}{Path=fonts/,Extension=.ttf,UprightFont=*-Regular,BoldFont=*-Bold,ItalicFont=*-Italic,BoldItalicFont=*-BoldItalic,Scale=0.9}
\setmainfont{LiberationSerif}[\optserif]\newfontfamily\cyrillicfont{LiberationSerif}[\optserif]
\setsansfont{Inter}[\optsans]\newfontfamily\cyrillicfontsf{Inter}[\optsans]
\setmonofont{LiberationMono}[\optmono]\newfontfamily\cyrillicfonttt{LiberationMono}[\optmono]
\usepackage[a4paper,landscape,left=1.8cm,right=1.8cm,top=1.7cm,bottom=1.9cm,columnsep=1.0cm]{geometry}
\usepackage{graphicx,xcolor,amsmath,array,booktabs,enumitem,titlesec,fancyhdr,caption}
\usepackage[most]{tcolorbox}
\usepackage[hidelinks]{hyperref}
\definecolor{cB}{HTML}{2A78D6}\definecolor{cO}{HTML}{C9501F}\definecolor{cA}{HTML}{12875C}\definecolor{cY}{HTML}{A87200}\definecolor{cM}{HTML}{52514E}
\definecolor{pA}{HTML}{EAF7F1}\definecolor{pY}{HTML}{FCF4DD}\definecolor{pB}{HTML}{EDF3FB}
\raggedright\setlength{\parindent}{0pt}\setlength{\parskip}{5pt}\emergencystretch=2em
\titleformat{\section}{\sffamily\Large\bfseries\color{cB}}{\thesection.}{0.5em}{}
\titleformat{\subsection}{\sffamily\large\bfseries}{}{0em}{}
\titlespacing*{\section}{0pt}{16pt}{6pt}\titlespacing*{\subsection}{0pt}{8pt}{2pt}
\setlist[itemize]{leftmargin=1.3em,itemsep=2pt,topsep=2pt,parsep=0pt}
\setlist[enumerate]{leftmargin=1.6em,itemsep=3pt,topsep=2pt,parsep=0pt}
\pagestyle{fancy}\fancyhf{}\renewcommand{\headrulewidth}{0pt}\fancyfoot[C]{\sffamily\small\color{cM}\thepage}
\captionsetup{font={small,sf},labelfont=bf,labelsep=period,justification=raggedright,singlelinecheck=false,skip=4pt}
\newcommand{\tU}{\textsuperscript{\sffamily\bfseries\color{cM}У}}
\newcommand{\tS}{\textsuperscript{\sffamily\bfseries\color{cB}С}}
\newcommand{\tP}{\textsuperscript{\sffamily\bfseries\color{cA}П}}
\newcommand{\tD}{\textsuperscript{\sffamily\bfseries\color{cO}Д}}
\newtcolorbox{note}{colback=pA,colframe=pA,boxrule=0pt,arc=2pt,left=8pt,right=8pt,top=5pt,bottom=5pt,before skip=8pt,after skip=8pt,fontupper=\small,title={\sffamily\bfseries\small\color{cA}Осторожно},coltitle=cA,colbacktitle=pA,attach title to upper={\par\vspace{1pt}}}
\newtcolorbox{ineq}{colback=pY,colframe=pY,boxrule=0pt,arc=2pt,left=8pt,right=8pt,top=5pt,bottom=5pt,before skip=8pt,after skip=8pt,fontupper=\small,title={\sffamily\bfseries\small\color{cY}Для математика},coltitle=cY,colbacktitle=pY,attach title to upper={\par\vspace{1pt}}}
\newcommand{\fig}[3][\textwidth]{\begin{figure*}[t]\centering\includegraphics[width=#1]{#2}\caption{#3}\end{figure*}}
\newcommand{\figp}[3][\textwidth]{\begin{figure*}[p]\centering\includegraphics[width=#1]{#2}\caption{#3}\end{figure*}}
\renewcommand{\dbltopfraction}{0.9}\renewcommand{\dblfloatpagefraction}{0.75}\renewcommand{\textfraction}{0.05}\renewcommand{\topfraction}{0.9}\renewcommand{\floatpagefraction}{0.75}\setcounter{dbltopnumber}{2}\setcounter{totalnumber}{4}\setlength{\dbltextfloatsep}{14pt}
\newcommand{\Hp}{H\textsubscript{2}O\textsubscript{2}}
\newcolumntype{R}[1]{>{\raggedright\arraybackslash}p{#1\linewidth}}

\newcommand{\facts}{\par\vspace{3pt}{\sffamily\bfseries\small\color{cM}Опорные факты}\par\vspace{-2pt}}

\raggedbottom
\begin{document}
\twocolumn[{%
{\sffamily\bfseries\fontsize{24}{28}\selectfont\color{cB} О шизофрении для математика\par}
\vspace{4pt}
{\sffamily\large Общая картина процессов: от молекулы до симптома\par}
\vspace{2pt}
{\sffamily\small\color{cM} Вводная брошюра к проекту о модели повреждения тормозных клеток коры. Версия 4.0 от 15 октября 2026 г.\par}
\vspace{14pt}}]

Брошюра написана для читателя, который владеет динамическими системами, но не занимался биологией и не читал статей по теме. Её цель — дать связную картину: какие процессы идут в клетке и в сети клеток, на каких временах, и где в этой картине появляется нелинейность, из-за которой болезнь может начинаться скачком. Подробностей мало намеренно: каждое утверждение дано на уровне, достаточном, чтобы понять следующее.

Разделы устроены одинаково. Сначала принцип простыми словами. Затем короткий список опорных фактов с метками происхождения. В~конце — рамка «Для математика»: каким математическим объектом это описывается.

\subsection{Как читать метки}
\begin{tabular}{@{}l@{\ \ }p{0.9\linewidth}@{}}
\tU & учебное знание: общепринятое, записано по~памяти и~отдельно не~сверялось;\\
\tS & взято из~статьи, прочитанной авторами проекта (числа — по~рабочим записям);\\
\tP & сверено поиском при подготовке брошюры, по~аннотации, препринту или пересказу;\\
\tD & допущение или расчёт авторов проекта, в~литературе в~таком виде нет.\\
\end{tabular}

Все рисунки — собственные схемы, нарисованные для брошюры.

\subsection{Содержание}
{\small
\begin{tabular}{@{}r@{\ \ }l@{\qquad}r@{\ \ }l@{}}
1 & Вся картина на одной странице & 14 & Соседняя гипотеза: BDNF\\
2 & Болезнь: что нужно объяснить & 15 & Откуда это известно\\
3 & Масштабы времени & 16 & Свидетельства: люди и животные\\
4 & Клетка за пять минут & 17 & Как из этого получается модель\\
5 & Нейрон и синапс & 18 & Переключатель сна\\
6 & Кора: возбуждение и торможение & 19 & От одной клетки к популяции\\
7 & NMDA-рецептор & 20 & Кора в петле\\
8 & Окислители и защита & 21 & Закон сохранения разрывов\\
9 & Переключатель & 22 & Что модель уже не смогла\\
10 & Петли и окно развития & 23 & Дополнительные результаты\\
11 & Сутки и сон & 24 & Открытые вопросы\\
12 & Шум & 25 & Словарь\\
13 & Дофамин и психоз & 26 & Медленная защита\\
 & & 27 & Прайминг\\
 & & 28 & Предвестник: кто его делает\\
 & & 29 & Из чего состоит эпизод\\
 & & 30 & Сколько клеток в популяции\\
 & & 31 & Как ломается ритм\\
 & & 32 & Как мы теперь считаем\\
 & & 33 & Открытые вопросы, словарь, источники\\
\end{tabular}}

Разделы 18–23 добавлены в~версии~3.0: они описывают, как модель устроена теперь и~что она уже успела запретить. Разделы 24–30 добавлены в~версии~4.0: медленная защита и~прайминг, приписка предвестника, устройство эпизода, число клеток как параметр, поломка ритма и~то, как устроен счёт. Первые семнадцать разделов — биология, они почти не~менялись.

% =====================================================================
\section{Вся картина на одной странице}
\fig[0.8\textwidth]{fig_overview.png}{Рабочая гипотеза одной цепочкой. Серые стрелки — причина и следствие; оранжевые — обратные связи, которые замыкают цепочку в~круг.}
Гипотеза, вокруг которой построена брошюра, состоит из шести шагов (рис.~1).
\begin{enumerate}
\item В~коре мозга есть небольшая группа тормозных клеток — \emph{PV-клетки}. Они удерживают возбуждение остальных клеток в~рамках и~задают ритм их работы.
\item PV-клетки работают на~пределе: много входов, высокая частота, большой расход энергии. Побочный продукт расхода энергии — окислители, в~первую очередь перекись водорода.
\item От~окислителей клетку защищают ферменты. Часть этой защиты клетка производит только пока исправно работает один из~её входных рецепторов — \emph{NMDA-рецептор}.
\item Если рецептор работает слабо (из-за наследственности или среды), защиты меньше, окислителей больше. Окислители, в~свою очередь, подавляют сам рецептор. Получается круг, который усиливает сам себя.
\item Перегруженная PV-клетка не~гибнет, а~«замолкает»: производит меньше своих рабочих белков. Торможение слабеет, кора и~гиппокамп перевозбуждаются; это добавляет окислителей — второй круг.
\item Перевозбуждённый гиппокамп через цепочку промежуточных отделов повышает чувствительность дофаминовой системы. Избыток дофамина связывают с~главными симптомами психоза — галлюцинациями и~бредом.
\end{enumerate}
\begin{note}
Каждая стрелка этой цепочки показана в~опытах по~отдельности: в~культуре клеток, на~срезах мозга или на~животных. Вся цепочка целиком у~человека не~прослежена.\par\vspace{4pt}
Изложена одна гипотеза — слабая работа NMDA-рецептора вместе с~нарушением окислительной защиты PV-клеток. Её источники: обзор Hardingham и~Do (2016), обзоры группы Cuenod и~Do (2021–2022), дофаминовая схема Грейса.\tS\ Это рабочая гипотеза, а~не~установленная картина болезни. Другие гипотезы ставят на~первое место сам дофамин, избыточное удаление синапсов в~подростковом возрасте, воспаление\tU\ или недостаток фактора роста BDNF и~связанные с~ним изменения серотониновой системы (раздел~14).
\end{note}
\begin{ineq}
Круг, усиливающий сам себя, — положительная обратная связь. Если её усиление достаточно велико, у~системы появляются два устойчивых состояния и~порог между ними. Тогда медленное накопление причин даёт резкий переход, а~возврат требует большего, чем потребовалось для срыва. Такой формой течения болезнь и~отличается (раздел~2). Насколько велико усиление на~самом деле — главный количественный вопрос проекта.
\end{ineq}

% =====================================================================
\section{Болезнь: что нужно объяснить}
Шизофрения — хроническое психическое расстройство. Международная классификация болезней последнего пересмотра (МКБ-11; действует с~2022~года, страны переходят на~неё постепенно) описывает её не~как набор готовых форм, а~как профиль по~нескольким независимым осям симптомов. Для модели важны две вещи: этот профиль у~разных больных разный, а~течение имеет характерную форму — долгая скрытая подготовка, сравнительно резкий переход и~затем эпизоды.
\facts
\begin{itemize}
\item Диагноз по~МКБ-11 требует не~меньше двух групп симптомов из~семи, присутствующих большую часть времени не~меньше месяца. Хотя~бы одна должна быть из~четырёх основных: бред (устойчивые ложные убеждения); галлюцинации, чаще всего голоса; расстройство мышления; переживание воздействия или управления извне. Остальные три: негативные симптомы (стало меньше эмоций, побуждений, речи), расстройство поведения, двигательные нарушения.\tP
\item Прежние формы (параноидная, гебефреническая, кататоническая, простая) в~МКБ-11 отменены. Вместо них у~каждого больного отдельно оценивают шесть осей: позитивные, негативные, депрессивные, маниакальные, двигательные и~когнитивные симптомы (память, внимание), а~также тип течения: первый эпизод, повторные эпизоды, непрерывное.\tP
\item Кататония — состояние с~резкими двигательными нарушениями — выделена в~отдельный диагноз.\tP
\item Рядом с~шизофренией в~той~же группе «первичных психотических расстройств» стоят отдельные диагнозы, при которых часть картины отсутствует или она кратковременна: шизоаффективное, шизотипическое, бредовое, острое преходящее психотическое расстройство.\tP
\item \emph{Психоз} — состояние с~выраженными позитивными симптомами; он приходит эпизодами.\tU
\item Распространённость в~течение жизни — 4,0 на~1000 человек (медиана по~исследованиям); риск заболеть за~жизнь — 7,2 на~1000, то~есть примерно один человек из~ста сорока.\tP
\item Начало чаще всего приходится на~конец подросткового возраста и~молодость.\tU
\item Наследуемость около 79~\%.\tP\ Это доля разброса риска в~популяции, связанная с~генами, а~не~вероятность передать болезнь ребёнку. Одного «гена шизофрении» нет: известны сотни вариантов с~малым вкладом каждого.\tU
\item Перед первым эпизодом идёт \emph{продром} — месяцы и~годы неспецифических изменений. Распределение его длительности имеет тяжёлый хвост: в~одном исследовании среднее 17,6, а~медиана 4~месяца.\tS
\item У~43,7~\% больных первый психоз развился остро, за~дни.\tS
\item Лекарства (антипсихотики) ослабляют позитивные симптомы и~почти не~действуют на~негативные и~когнитивные.\tU
\end{itemize}
\begin{note}
Числа о~распространённости, наследуемости и~течении получены при прежних критериях диагноза (МКБ-9, МКБ-10 и~американские руководства); при переходе к~МКБ-11 границы диагноза сдвинулись. Долгая подготовка и~резкий переход наводят на~мысль о~системе с~порогом, но~постепенное накопление изменений без всякого порога дало~бы похожую картину, если симптомы появляются только после некоторого уровня повреждения. Форма течения — довод, а~не~доказательство.
\end{note}
\begin{ineq}
Диагноз — это точка в~пространстве шести осей, а~не~значение одной переменной. Модель одного механизма может претендовать на~объяснение отдельных осей, но~не~болезни целиком. Гипотеза этой брошюры относится прежде всего к~когнитивной оси (через PV-клетки и~ритмы коры) и~к~позитивной (через дофамин); о~депрессивной и~маниакальной она ничего не~говорит.
\end{ineq}

% =====================================================================
\section{Масштабы времени}
\figp[0.9\textwidth]{fig_times.png}{Характерные времена явлений, о~которых идёт речь в~брошюре. Шкала логарифмическая.}
Процессы, имеющие отношение к~болезни, занимают больше десяти порядков по~времени: от~миллисекундного импульса нейрона до~многолетнего созревания коры (рис.~2). Быстрые процессы для медленных выглядят как мгновенно установившиеся; медленные для быстрых — как постоянные параметры.
\facts
\begin{itemize}
\item Импульс нейрона длится около миллисекунды.\tU\ Галлюцинация — обычно секунды; только 7,4~\% эпизодов продолжаются 10~минут и~дольше.\tS
\item Количество белка в~клетке меняется за~часы и~сутки.\tU
\item Способность мозга синтезировать дофамин растёт месяцы и~годы — от~продрома к~психозу.\tS
\item Внешние ритмы задают свои масштабы: сутки (свет, сон и~бодрствование) и~год (длина дня, температура).\tU
\end{itemize}
\begin{ineq}
Разделение времён позволяет строить модель слоями: быстрые переменные исключаются как алгебраические связи, медленные замораживаются как параметры. Одна-две медленные переменные такого разброса не~покроют; модель проекта описывает середину шкалы — от~минут до~суток, события внутри одной клетки.
\end{ineq}

% =====================================================================
\section{Клетка за пять минут}
Клетка — капля водного раствора в~оболочке (\emph{мембране}). Почти всю работу в~ней делают \emph{белки}. Одни ускоряют химические реакции (их называют \emph{ферментами}), другие служат воротами в~мембране, третьи принимают сигналы извне (\emph{рецепторы}).

Белок строится по~гену в~два шага: с~гена снимается рабочая копия (мРНК), по~ней собирается белок. Всё вместе называют \emph{экспрессией} гена. Готовый белок живёт ограниченное время и~разрушается. Поэтому количество любого белка — это баланс синтеза и~распада, а~клетка управляет собой, меняя скорость синтеза: «включает» и~«выключает» гены в~ответ на~сигналы.

Энергию клетка получает в~\emph{митохондриях}, где питательные вещества окисляются кислородом. Энергия запасается в~молекулах АТФ и~тратится главным образом на~поддержание разности концентраций ионов по~разные стороны мембраны.
\facts
\begin{itemize}
\item Время жизни разных белков — от~минут до~недель.\tU
\item Скорость ферментативной реакции растёт с~концентрацией исходного вещества и~выходит на~насыщение: фермента ограниченное количество, и~каждая его молекула занята конечное время.\tU
\item Сигнал от~рецептора до~гена идёт по~цепочке промежуточных молекул; от~сигнала до~заметного изменения количества белка проходят часы.\tU
\item Большинство величин в~клетке удерживается отрицательными обратными связями.\tU
\item Мозг составляет около 2~\% массы тела и~потребляет около 20~\% кислорода.\tU
\end{itemize}
\begin{ineq}
Количество белка $p$ подчиняется уравнению $p' = s(t) - p/\tau$: это звено первого порядка с~постоянной времени, равной времени жизни белка. Ферментативная реакция — насыщающаяся нелинейность вида $v = V s/(K + s)$. Управление генами даёт медленные переменные, химия малых молекул — быстрые. Насыщение защиты и~есть источник нелинейности, на~которой построен раздел~9.
\end{ineq}

% =====================================================================
\section{Нейрон и синапс}
\fig{fig_neuron.png}{Нейрон и~синапс. Слева — клетка целиком, справа — увеличенное место контакта двух клеток.}
Нейрон — клетка, приспособленная для передачи сигналов. На~мембране он поддерживает электрическое напряжение. Когда суммарный входной сигнал превышает порог, по~выходному отростку пробегает короткий импульс. Между клетками сигнал передаётся химически: импульс выбрасывает из~окончания молекулы \emph{медиатора}, а~они открывают рецепторы на~следующей клетке (рис.~3).
\facts
\begin{itemize}
\item Нейрон состоит из~тела, ветвящихся \emph{дендритов} (вход) и~одного \emph{аксона} (выход). Импульс устроен по~принципу «всё или ничего»; информацию несут частота и~время импульсов.\tU
\item Главный возбуждающий медиатор коры — \emph{глутамат}, главный тормозной — \emph{ГАМК}. Возбуждение приближает клетку к~порогу, торможение отдаляет.\tU
\item У~глутамата два основных рецептора. \emph{AMPA} — быстрый, срабатывает за~миллисекунды. \emph{NMDA} — медленный (десятки и~сотни миллисекунд), о~нём раздел~7.\tU
\item Кальций, вошедший в~клетку при её работе, служит сигналом для включения генов. Так активность синапсов управляет тем, какие белки клетка производит.\tU
\end{itemize}
\begin{ineq}
Нейрон — пороговый элемент, синапс — связь со~знаком и~весом. Важнее для дальнейшего другое: через кальций быстрая электрическая активность (миллисекунды) управляет медленной химией клетки (часы). Это мост между масштабами времени.
\end{ineq}

% =====================================================================
\section{Кора: возбуждение и торможение}
\fig{fig_circuit.png}{Простейшая схема коры: возбуждающие (пирамидные) клетки возбуждают тормозную клетку, а~она тормозит их в~ответ. Справа — три молекулы, о~которых идёт речь дальше.}
Большинство нейронов коры возбуждают друг друга. Такая сеть сама по~себе неустойчива: активность разгонялась~бы до~предела. Её удерживают тормозные клетки, которые получают возбуждение от~соседей и~тормозят их в~ответ (рис.~4). Самые быстрые из~них содержат белок \emph{парвальбумин} и~называются PV-клетками.
\facts
\begin{itemize}
\item В~коре около 80~\% нейронов возбуждающие (\emph{пирамидные}) и~около 20~\% тормозные.\tU\ PV-клетки составляют примерно 40~\% тормозных клеток коры мыши.\tU
\item Одна PV-клетка получает 16–34~тысячи синапсов; торможение приходит к~пирамидной клетке с~задержкой не~больше 2~мс; частота разрядов достигает сотни импульсов в~секунду.\tS
\item У~PV-клеток высокий расход энергии.\tS
\item Парвальбумин связывает кальций и~быстро убирает его после импульса.\tU\ \emph{GAD67} — фермент, который делает ГАМК: меньше GAD67 — слабее торможение.\tU
\item Вокруг PV-клеток есть оболочка из~внеклеточного вещества (\emph{перинейрональная сеть}). Она защищает клетку от~окислителей и~созревает в~подростковом возрасте.\tS
\item PV-клетки задают \emph{гамма-ритм} — колебания активности коры с~частотой 30–80~Гц. У~больных шизофренией этот ритм нарушен.\tU
\item Кора работает в~режиме, где возбуждение само себя усиливает и~удерживается только торможением.\tS
\end{itemize}
\begin{ineq}
PV-клетка — быстрая отрицательная обратная связь, без которой сеть неустойчива, и~одновременно её тактовый генератор. Плата за~скорость — расход энергии, а~значит, и~поток окислителей (раздел~8). Это самое нагруженное звено схемы.
\end{ineq}

% =====================================================================
\section{NMDA-рецептор}
NMDA-рецептор открывается только при совпадении двух условий: медиатор пришёл \emph{и}~клетка уже возбуждена. Через него в~клетку входит кальций. Поэтому он служит не~только входом, но~и~датчиком активности: по~его работе клетка «узнаёт», что она нужна сети, и~поддерживает производство своих белков — и~рабочих, и~защитных.
\facts
\begin{itemize}
\item Кетамин блокирует NMDA-рецептор. У~здоровых добровольцев он вызывает состояния, похожие на~симптомы шизофрении, и~ухудшает внимание и~память.\tP
\item Парадокс: блокируют возбуждающий рецептор, а~активность коры растёт. Принятое объяснение: сильнее страдают рецепторы на~тормозных клетках, торможение слабеет, пирамидные клетки \emph{растормаживаются}.\tU
\item В~культуре клеток кетамин снижает количество GAD67 вдвое примерно за~4~часа, а~парвальбумина — за~8~часов. Эффект обратим, клетки не~гибнут.\tS
\item У~мышей одна доза кетамина следа не~оставляет, а~две дозы в~два дня подряд вызывают стойкую потерю парвальбумина и~GAD67 с~медленным и~неполным восстановлением к~10-му дню.\tS
\item Блокада рецептора у~новорождённых грызунов за~сутки вдвое снижает активность фермента, производящего главный антиоксидант клетки.\tS\ Синаптическая активность включает в~нейронах гены защитных ферментов, в~том числе того, который чинит повреждённую защиту (раздел~9).\tP
\item Окислители обратимо подавляют сам NMDA-рецептор.\tS
\end{itemize}
\begin{ineq}
Три факта складываются в~замкнутую петлю: рецептор поддерживает защиту; без защиты растут окислители; окислители подавляют рецептор. Опыт с~двумя дозами показывает порог не~только по~силе, но~и~по~\emph{длительности} воздействия: у~системы есть память длиной хотя~бы в~сутки.
\end{ineq}

% =====================================================================
\section{Окислители и защита}\label{sec:redox}
\figp{fig_redox.png}{Откуда берётся перекись и~куда уходит. Внизу — цикл защитного фермента пероксиредоксина с~боковой ветвью, которая его выключает.}
Окислитель — молекула, отнимающая электроны у~других молекул. Отняв электрон у~белка или у~ДНК, она их портит. Окислители неизбежно образуются там, где клетка работает с~кислородом. Главный из~них для дальнейшего — перекись водорода \Hp: она сравнительно долго живёт и~проходит сквозь мембраны. В~малых количествах перекись — обычный сигнал внутри клетки, в~больших — яд. \emph{Окислительным стрессом} называют состояние, когда приток окислителей превышает возможности очистки.

Защита устроена в~два слоя. Первый — ферменты, которые разлагают перекись; самые быстрые из~них — \emph{пероксиредоксины}. Второй — запас малых молекул, которые отдают электроны и~этим возвращают ферменты в~рабочее состояние; главный такой запас — \emph{глутатион}.
\facts
\begin{itemize}
\item Источники: митохондрии (побочный продукт выработки энергии)\tU\ и~фермент Nox2, который производит окислители намеренно. В~нейронах Nox2 включается при возбуждении NMDA-рецепторов; его продукт выходит наружу и~действует на~соседние клетки.\tS
\item Обычная концентрация перекиси в~клетке — 1–10~нМ.\tS
\item Глутатион — короткая молекула из~трёх аминокислот; его в~клетке на~пять-шесть порядков больше, чем перекиси.\tU\ Половина глутатиона мозга мыши обновляется за~60–80~часов.\tP\ При этом его уровень в~мозге мыши меняется в~течение суток в~1,2~раза.\tS
\item Сырьё для глутатиона нейроны получают от~\emph{астроцитов} — опорных клеток мозга.\tU
\item В~зрелых нейронах система, включающая защитные гены в~ответ на~сам стресс, намного слабее, чем в~астроцитах.\tS
\end{itemize}
\begin{ineq}
Баланс перекиси быстрый: приток равен оттоку за~секунды, так что концентрация перекиси — алгебраическая функция притока и~состояния защиты. Нелинейность возникает оттого, что отток насыщается и, хуже того, сама защита портится от~избытка перекиси (следующий раздел).
\end{ineq}

% =====================================================================
\section{Переключатель}
\fig{fig_switch.png}{Слева: установившаяся доля выключенного фермента в~зависимости от~притока перекиси (расчёт по~упрощённому уравнению, параметры подобраны для наглядности). Справа: та~же картина в~виде рельефа.}
Пероксиредоксин обезвреживает молекулу перекиси, окисляясь сам, и~за~секунды возвращается в~рабочее состояние. Но~если к~уже окисленному ферменту успевает прийти вторая молекула перекиси, он \emph{переокисляется} и~выключается надолго. Чинит его отдельный фермент, \emph{сульфиредоксин}, и~делает это медленно.

Отсюда положительная обратная связь: чем больше перекиси, тем больше фермента выключено; чем меньше рабочего фермента, тем хуже очистка и~тем больше перекиси. При достаточно сильной связи у~клетки появляются два устойчивых состояния при одной и~той~же нагрузке: «защита работает» и~«защита выключена» (рис.~6).
\facts
\begin{itemize}
\item Починка переокисленного фермента занимает от~3~минут до~часа в~клетках человека.\tS
\item В~митохондриях починка идёт иначе: сульфиредоксина там в~покое почти нет. Он ввозится после того, как перекись выходит из~митохондрий наружу, достигает максимума через 4–8~часов и~разрушается, когда починка закончена.\tS
\item Подробная модель этой системы (Selvaggio и~соавт., 2018; параметры измерены для 15~типов клеток) даёт два устойчивых состояния только у~одной клеточной линии и~в~узком диапазоне нагрузок.\tS
\item Упрощение этой модели до~одного уравнения даёт точный критерий двух состояний. Измеренные клетки лежат у~самой границы, но~снаружи: чтобы попасть внутрь, починку нужно ослабить в~1,6–22~раза.\tD
\item Порог гибели клетки резкий: удвоение дозы перекиси переводит клетки от~полной сохранности к~почти полной гибели.\tS
\end{itemize}
\begin{note}
Обычные клетки \emph{не}~находятся в~области двух состояний. Чтобы переключатель заработал, защита должна быть ослаблена. Гипотеза утверждает, что у~PV-клеток при слабом NMDA-рецепторе так и~есть; прямых измерений для этих клеток нет. Кинетика фермента измерена в~клеточных линиях, не~в~нейронах.
\end{note}
\begin{ineq}
Это пара складок (седло-узловых бифуркаций) с~гистерезисом: срыв и~возврат происходят при разных нагрузках. Внутри окна переход возможен только под действием возмущения — толчка или шума (раздел~12). Ширина окна и~высота барьера — те~величины, которые нужно оценивать количественно.
\end{ineq}

% =====================================================================
\section{Петли и окно развития}
\fig[0.7\textwidth]{fig_loops.png}{Петли положительной обратной связи, описанные в~обзорах, и~место подросткового возраста на~шкале жизни.}
Петля из~раздела~7 — не~единственная. Обзоры описывают несколько кругов, проходящих через окислительный стресс, и~все они сходятся на~PV-клетке (рис.~7). Отдельный вопрос — почему болезнь начинается в~молодости. Ответ биологов: PV-клетки и~их защитные оболочки созревают поздно, к~концу подросткового возраста, и~до~этого уязвимы.
\facts
\begin{itemize}
\item Описаны две основные петли между слабой работой NMDA-рецептора и~окислительным стрессом: внутриклеточная (рецептор~→ защитные гены) и~сетевая (растормаживание~→ сигнал воспаления~→ Nox2~→ окислители).\tS
\item Другой обзор насчитывает пять петель: через рецептор, митохондрии, иммунные клетки мозга, оболочки вокруг клеток и~дофамин. Чисел в~обзоре нет.\tS
\item Генетический риск обычно компенсируется; исход зависит от~силы \emph{и~длительности} стресса; мышам с~выключенной частью рецептора для нарушений нужен «второй удар».\tS
\item Окислительный стресс повреждает PV-клетки именно в~детском и~подростковом возрасте.\tS
\item У~крыс короткий курс успокоительного в~подростковом возрасте предотвращает последующие нарушения.\tS
\end{itemize}
\begin{ineq}
«Порочный круг» в~биологическом обзоре — это стрелки на~рисунке. Замкнётся~ли петля, зависит от~произведения усилений вдоль неё: если оно меньше единицы, возмущение затухает. Обзоры этого вопроса не~ставят, потому что у~них нет чисел. Окно развития — медленный дрейф параметров защиты, то~есть медленный проход через границу области двух состояний.
\end{ineq}

% =====================================================================
\section{Сутки и сон}
\fig{fig_day.png}{Что и~когда достигает максимума и~минимума за~сутки у~грызунов. Данные из~разных работ, видов и~отделов мозга; для нижней строки время указано приблизительно.}
Почти в~каждой клетке тела идут собственные часы с~периодом около суток. Их механизм — петля отрицательной обратной связи с~запаздыванием: два белка включают гены, продукты которых через несколько часов подавляют эти~же два белка. Главные часы находятся в~небольшом отделе мозга над перекрёстом зрительных нервов; свет через глаза подстраивает их под смену дня и~ночи, а~они синхронизируют остальные.

Для нашей темы сутки важны по~двум причинам. Во-первых, сон и~бодрствование меняют активность нейронов, а~с~ней и~приток окислителей: нагрузка на~защиту не~постоянна, она каждый день проходит цикл. Во-вторых, сама окислительная защита колеблется с~суточным периодом, причём эти колебания найдены даже в~клетках без ядра, то~есть без генов.
\facts
\begin{itemize}
\item Собственный период часов сохраняется в~постоянной темноте.\tU\ Суточный ритм переокисления пероксиредоксина есть в~красных клетках крови, у~которых нет ядра.\tU
\item В~главных часах мозга мыши количество переокисленного пероксиредоксина и~общее количество этого фермента меняются за~сутки в~несколько раз; переокисленная форма опережает общее количество примерно на~4~часа (чтение графика; три животных).\tS
\item В~срезах главных часов ткань восстановлена днём и~окислена в~начале ночи; изменение окислительного состояния меняет возбудимость нейронов примерно за~минуту.\tS
\item В~PV-клетках лобной коры крысы метка окислительного повреждения ДНК максимальна к~концу активной фазы и~к~концу фазы покоя падает на~47~\% (рис.~8).\tP
\item Шесть часов нарушенного сна повышают эту метку у~крыс на~25~\%; в~PV-клетках она на~27–50~\% выше, чем в~соседних.\tS\ После 4~часов без сна частота разрядов нейронов коры выше на~34–64~\%.\tS
\item В~PV-клетках лобной коры мыши в~суточном ритме считывается 9~\% генов у~самцов и~14~\% у~самок (в~пирамидных клетках около 35~\%); у~самцов это в~основном гены белков митохондрий. Собственная возбудимость PV-клеток днём и~ночью одинакова.\tP
\item В~надпочечнике, сердце и~буром жире мыши переокисленный пероксиредоксин митохондрий и~чинящий его сульфиредоксин колеблются с~суточным периодом почти в~противофазе. Повреждение максимально к~концу активной фазы, сульфиредоксин — через 8–12~часов (чтение рисунка; шесть точек за~сутки). Общее количество самого пероксиредоксина при этом постоянно.\tS
\item Нарушения сна и~суточного ритма встречаются у~значительной части больных шизофренией.\tU
\end{itemize}
\begin{note}
Грызуны — ночные животные: их активная фаза приходится на~темноту. Переносить часы суток с~мыши на~человека напрямую нельзя; переносится фаза относительно сна и~бодрствования. Суточный ритм переокисленного пероксиредоксина в~коре и~в~PV-клетках не~измерен: данные есть для главных часов и~для органов вне мозга. В~пробе целого мозга мыши переокисленный пероксиредоксин митохондрий в~обычных условиях не~обнаружен;\tS\ о~редком типе клеток такая проба ничего не~говорит.
\end{note}
\begin{ineq}
Здесь два разных объекта. Первый — периодическое внешнее воздействие на~систему с~двумя состояниями: если суточный размах нагрузки сравним с~шириной окна, система каждый день подходит к~границе. Второй — собственный релаксационный цикл «повреждение → включение починки → восстановление → спад починки»; такая модель для клеток вне мозга опубликована (del~Olmo и~соавт., 2019).\tS\ Различить их можно по~сдвигу фаз: звено первого порядка не~может отстать от~своего входа больше чем на~четверть периода, а~наблюдаемые 8–12~часов из~24 — это от~трети до~половины.\tD
\end{ineq}

% =====================================================================
\section{Шум}
Ни~одна величина в~мозге не~стоит на~месте. Активность нейронов колеблется сама по~себе, без внешних причин, причём на~всех временах сразу: от~долей секунды до~минут и~дольше. Вместе с~ней колеблется приток окислителей. Для системы с~двумя состояниями это существенно: случайный толчок может перебросить её через порог раньше, чем до~порога дойдёт средняя нагрузка.
\facts
\begin{itemize}
\item Спектр самопроизвольной активности мозга не~имеет выделенного масштаба: мощность убывает с~частотой по~степенному закону в~диапазоне примерно от~0,003 до~500~Гц.\tS
\item Уровень бодрствования меняется скачками за~секунды и~возвращается к~прежнему за~десятки секунд.\tS
\item Эпизоды «голосов» у~больных возникают спонтанно, без внешнего повода.\tS
\end{itemize}
\begin{ineq}
Это задача о~выходе из~потенциальной ямы под действием шума. Среднее время ожидания перехода экспоненциально зависит от~отношения высоты барьера к~интенсивности шума. У~границы окна барьер убывает как расстояние до~границы в~степени $3/2$.\tD\ Экспоненциальная зависимость означает, что небольшой разброс параметров между клетками или людьми даёт огромный разброс времён — распределение с~тяжёлым хвостом, как у~длительности продрома. Расчёт проекта показал и~ограничение: быстрый шум правдоподобной силы слишком слаб, чтобы перебросить систему вдали от~границы; существенны медленные колебания, на~минутах и~дольше.\tD
\end{ineq}

% =====================================================================
\section{Дофамин и психоз}
\fig{fig_dopamine.png}{Путь от~гиппокампа к~дофаминовым клеткам по~работам группы Грейса. Стрелка — возбуждение, черта на~конце — торможение.}
Повреждение PV-клеток само по~себе не~объясняет, почему возникают голоса и~бред. Мостом к~симптомам считают \emph{дофамин} — медиатор, которым мозг помечает события как значимые. Если его выброс избыточен и~беспорядочен, значимыми кажутся случайные события, а~бред возникает как попытка их объяснить.
\facts
\begin{itemize}
\item Все применяемые антипсихотики закрывают рецепторы дофамина одного типа (D2). Эффект появляется, когда занято около 65~\% рецепторов; выше 78~\% начинаются двигательные побочные эффекты.\tP
\item По~сводке 44~исследований выработка и~выброс дофамина у~больных повышены, а~число рецепторов почти не~изменено.\tP
\item Вещества, усиливающие выброс дофамина, вызывают или обостряют психоз.\tU
\item Схема Грейса: ослабление PV-клеток в~гиппокампе делает его перевозбуждённым; через два тормозных звена это увеличивает долю дофаминовых клеток, готовых ответить на~сигнал (рис.~9).\tS
\end{itemize}
\begin{ineq}
Повреждение тормозных клеток меняет коэффициент усиления дофаминовой системы. В~модели проекта этого слоя нет: она заканчивается на~состоянии одной клетки. Симптомы — отдельная будущая задача, со~своим, секундным масштабом времени.
\end{ineq}

% =====================================================================
\section{Соседняя гипотеза: BDNF и серотонин}
Нейрону недостаточно получать питание: чтобы сохранять отростки и~синапсы, ему нужны сигналы поддержки от~соседей — \emph{нейротрофические факторы}. Главный из~них в~коре — белок BDNF. Его выделяют работающие возбуждающие клетки, а~принимает рецептор TrkB на~соседних клетках, в~том числе на~PV-клетках. \emph{Нейротрофическая гипотеза} шизофрении утверждает, что этой поддержки не~хватает: в~развитии это даёт неправильную сборку сетей, у~взрослого — сниженную способность сетей перестраиваться.

\emph{Серотонин} — ещё один медиатор, наряду с~дофамином управляющий общим режимом работы коры. Его связывают с~психозом давно: на~его рецепторы действуют галлюциногены. Один из~этих рецепторов управляет выработкой BDNF, так что две линии сходятся.
\facts
\begin{itemize}
\item В~лобной коре умерших больных количество мРНК BDNF ниже на~28~\%, а~его рецептора TrkB — на~27~\%. Чем сильнее у~больного снижен TrkB, тем сильнее снижен GAD67 (корреляция 0,74).\tP
\item У~мышей со~сниженным на~75~\% рецептором TrkB количество мРНК GAD67 и~парвальбумина ниже на~23 и~40~\%. У~мышей со~сниженным на~80~\% самим BDNF они не~изменены. Вывод авторов: дело в~сигнале через рецептор.\tP
\item В~крови больных BDNF умеренно снижен (сводка 41~исследования, более 7000~человек). Снижение растёт с~длительностью болезни и~не~связано с~тяжестью симптомов; причина это или следствие — авторы сводки оставляют открытым.\tP
\item Выработка BDNF зависит от~активности клетки и~от~кальция, входящего в~том числе через NMDA-рецептор.\tU
\item Серотониновый рецептор 5-HT2A управляет выработкой BDNF в~лобной коре и~гиппокампе; связь двусторонняя.\tP
\item Галлюциногены (ЛСД, псилоцибин) действуют через рецептор 5-HT2A; большинство антипсихотиков второго поколения закрывают не~только дофаминовые рецепторы D2, но~и~его.\tU
\item Антипсихотики и~антидепрессанты повышают выработку BDNF (по~обзору 2005~года).\tP
\end{itemize}
\begin{note}
Нейротрофическая гипотеза не~нова: обзоры под таким названием вышли в~2004–2005 годах. Она не~исключает гипотезу этой брошюры: обе описывают ослабление PV-клеток, но~ставят первым разное звено. Какое из~них первично и~не~являются~ли оба следствиями общей причины — не~установлено. Единой опубликованной модели, объединяющей серотонин, BDNF и~окислительный стресс, поиск при подготовке брошюры не~нашёл.
\end{note}
\begin{ineq}
С~точки зрения уравнений BDNF — ещё один кандидат на~роль медленной связи «активность сети → поддержка PV-клетки»: его выделяют активные клетки, а~действует он за~часы и~дни. Конструкция та~же, что в~разделе~7: медленная переменная, которую поддерживает активность и~которая сама поддерживает торможение. Отличие в~том, что источник поддержки — соседние возбуждающие клетки, то~есть петля замыкается через сеть, а~не~внутри одной клетки.\tD
\end{ineq}

% =====================================================================
\section{Откуда это известно}
Почти всё сказанное выше получено не~на~больных людях. Чем ближе способ измерения к~молекуле, тем дальше он от~болезни; надёжность вывода определяется тем, на~какой ступени он получен (таблица~1).
\begin{table*}[t]\small\centering
\begin{tabular}{@{}>{\raggedright\arraybackslash}p{0.17\textwidth}>{\raggedright\arraybackslash}p{0.37\textwidth}>{\raggedright\arraybackslash}p{0.38\textwidth}@{}}
\toprule
\textsf{\textbf{Ступень}} & \textsf{\textbf{Что позволяет}} & \textsf{\textbf{Чем приходится платить}}\\
\midrule
Очищенные белки и~клеточные линии & измерить скорости реакций; отсюда кинетика пероксиредоксина & это не~нейроны; условия далеки от~живой клетки\\
\addlinespace
Культура нейронов & часами следить за~клетками, добавлять вещества в~точной дозе & клетки взяты у~зародышей или новорождённых, они незрелые и~лишены обычного окружения\\
\addlinespace
Срез мозга & записывать ток отдельной клетки в~сохранной местной сети & срез живёт часы; дальние связи обрезаны\\
\addlinespace
Живое животное & вмешиваться и~смотреть последствия через дни и~месяцы; выключать отдельные гены & мышь не~болеет шизофренией; воспроизводится звено, а~не~болезнь\\
\addlinespace
Человек: ткань после смерти & считать клетки и~белки в~нужном отделе & конечное состояние после десятилетий болезни и~лечения\\
\addlinespace
Человек: прижизненные снимки & измерять у~больных и~в~динамике & среднее по~кубику ткани объёмом в~кубические сантиметры\\
\bottomrule
\end{tabular}
\caption{Ступени свидетельств: от молекулы к человеку.}
\end{table*}
\facts
\begin{itemize}
\item \emph{Окрашивание антителами}: срез ткани обрабатывают антителом к~нужному белку с~цветной меткой и~считают клетки или яркость. Клетка, переставшая производить белок-метку, выглядит как исчезнувшая. Поэтому «потеря PV-клеток» может означать и~гибель, и~молчание.\tU
\item \emph{Разделение белков в~геле}: белки разгоняют по~размеру, метят антителом и~измеряют яркость полосы. Результат относительный, обычно в~долях от~контроля или от~максимума.\tU
\item \emph{Запись тока клетки}: тонкий электрод на~мембране одной клетки в~срезе; даёт точные числа для немногих клеток.\tU
\item \emph{Снимки с~радиоактивной меткой} (ПЭТ) показывают, например, способность синтезировать дофамин; \emph{магнитно-резонансная спектроскопия} — концентрацию вещества, в~частности глутатиона, усреднённую по~всем клеткам объёма.\tU
\item Группы в~опытах на~животных малы: три–десять животных на~точку — обычное дело.\tU
\end{itemize}
\begin{ineq}
Для построения модели из~этого следуют три правила. Число, измеренное на~одной ступени, нельзя без оговорки подставлять на~другой. Относительные единицы дают отношения и~сдвиги фаз, но~не~абсолютные значения. Среднее по~объёму ткани ничего не~говорит о~редком типе клеток: PV-клетки составляют несколько процентов нейронов.
\end{ineq}

% =====================================================================
\section{Свидетельства: люди и животные}
\subsection{Что видно у больных}
\begin{itemize}
\item В~лобной коре умерших больных снижено количество мРНК парвальбумина и~GAD67, но~плотность самих PV-клеток не~изменена: клетки на~месте, но~производят меньше этих белков.\tP
\item Глутатион. Первая работа (2000~год): в~спинномозговой жидкости на~27~\% меньше, в~лобной коре на~52~\% меньше.\tS\ Сводка 2021~года (589~больных и~601~здоровый): средняя разница незначима, но~разброс у~больных больше.\tS
\item Способность синтезировать дофамин повышена (раздел~13).\tP
\end{itemize}
\subsection{Что показано в опытах}
\begin{itemize}
\item \textbf{Кетамин в~культуре клеток.}\tS\ GAD67 и~парвальбумин падают вдвое за~4 и~8~часов; обратимо. Даёт масштаб времени медленного звена и~довод в~пользу «молчания», а~не~гибели.
\item \textbf{Кетамин мышам.}\tS\ Одна доза — ничего; два дня подряд — стойкая потеря. Порог по~длительности.
\item \textbf{Блокада рецептора у~новорождённых грызунов.}\tS\ Синтез глутатиона падает вдвое за~сутки. Связь «рецептор → защита».
\item \textbf{Повреждение мозга в~утробе у~крыс} (вещество MAM).\tS\ Потеря парвальбумина в~гиппокампе, затем избыток дофаминовой активности у~взрослых. Мост к~дофамину и~роль подросткового возраста.
\item \textbf{Мыши без части NMDA-рецептора.}\tS\ Для нарушений нужен второй удар: слабый рецептор сдвигает к~границе, но~не~переводит через неё.
\end{itemize}
\begin{note}
Отсутствие среднего дефицита глутатиона у~больных гипотезу не~опровергает: спектроскопия усредняет по~всем клеткам, а~больший разброс может означать подгруппу с~глубоким дефицитом или нарушенный суточный ритм. Но~и~не~подтверждает: объяснение «дело в~подгруппе или в~редком типе клеток» требует независимой проверки, иначе оно делает гипотезу неопровержимой.
\end{note}

% =====================================================================
\section{Как из этого получается модель}
\fig[0.86\textwidth]{fig_model.png}{Ядро модели. Острая стрелка — усиление, черта на~конце — подавление. Цвет показывает, насколько связь обоснована.}
Готовые расчётные модели есть для краёв картины. Быстрые последствия потери парвальбумина и~GAD67 для сети смоделированы (Volman и~соавт., 2011); кинетика пероксиредоксина — тоже (Selvaggio и~соавт., 2018), как и~её суточные колебания вне мозга (del~Olmo и~соавт., 2019). Модели самой петли повреждения — рецептор, окислители, тормозная клетка — авторы проекта в~литературе не~нашли.\tS

Ядро модели проекта — две переменные (рис.~10): быстрая $x$, доля рабочего пероксиредоксина, и~медленная $c$, уровень защиты, который поддерживается работой рецептора. Концентрация перекиси $u$ находится из~алгебраического баланса с~притоком $\nu$:
\begin{align*}
&\nu = \varepsilon u + \frac{x\,u}{1+u},\quad \varepsilon = e_0 c^{\gamma},\qquad N = \frac{N_0}{1 + u/K_N},\\
&x' = c^{m}(1 - x) - \mu_0\,\frac{x\,u^2}{1+u},\\
&c' = \frac{1}{T}\Big(N + \frac{\beta u}{K_B + u} - c\Big).
\end{align*}
Первая строка и~уравнение для $x$ выведены из~измеренной кинетики. Уравнение для $c$ — допущение о~форме: сама связь «рецептор → защита» описана в~статьях. Время измеряется в~единицах времени починки фермента; $T$ — отношение времени жизни защитных белков к~нему. Параметр $N_0 < 1$ изображает врождённо слабый рецептор.
\subsection{Что посчитано}
\begin{itemize}
\item Медленная связь через рецептор расширяет окно двух состояний: медиана отношения его границ растёт с~1,46 до~1,88; в~10,5~\% случайных наборов параметров два состояния появляются только благодаря ей.\tD
\item Отрицательный результат. Ядро не~даёт несимметрии «быстро теряется — долго возвращается» и~не~воспроизводит опыт с~двумя дозами кетамина: памяти между дозами в~нём нет. Добавление глутатиона и~обратной связи через активность сети этого не~исправило.\tD
\item В~4,6~\% наборов ядро даёт незатухающие колебания «повреждение — восстановление» с~периодом около 3,4~$T$. Повреждение опережает защиту на~0,16 периода — меньше четверти, как и~должно быть у~звена первого порядка; наблюдаемый для сульфиредоксина сдвиг — от~трети до~половины периода — ядру недоступен.\tD
\item Со~случайными колебаниями нагрузки закон «барьер убывает как расстояние до~границы в~степени $3/2$» и~формула для времени ожидания перехода подтверждены расчётом. Барьеры оказались высокими для быстрого шума.\tD
\end{itemize}

% =====================================================================
\section{Переключатель сна}\label{sec:sleep}
\fig{fig_sleepswitch.png}{Слева: два ядра тормозят друг друга, поэтому при одном и~том~же притоке возможны два состояния. Стрелки — переходы на~складках. Справа: накопленное за~день «давление сна» при одном и~том~же расписании.}

Нагрузка на~клетку не~постоянна: днём мозг работает, ночью отдыхает. Чтобы это вошло в~модель честно, расписание сна нельзя задавать руками — сон должен получаться сам.

Устроен он так. Есть «сонное» ядро и~есть ядра бодрствования; они тормозят друг друга. Такая пара ведёт себя как триггер: при одном и~том~же внешнем притоке у~неё два устойчивых состояния, и~система сидит в~том, в~котором оказалась. Переход происходит не~при пересечении какого-то порога, а~когда ветвь, на~которой система сидела, просто исчезает — такая точка называется складкой.

Притоком к~сонному ядру управляют две вещи: внутренние часы с~периодом около суток и~накопленные за~время бодрствования «вещества сна». Вторые растут, пока человек не~спит, и~убывают во~сне. Когда их становится достаточно, верхняя ветвь исчезает — человек засыпает; когда за~ночь их становится мало, исчезает нижняя — просыпается.

\facts
\begin{itemize}
\item Модель пары ядер и~все девять её чисел взяты из~статьи целиком и~не~настраивались. Проверка: складки у~нас вышли при 1,450 и~2,463~мВ, в~статье — 1,45 и~2,46; свободный цикл дал бодрствование 15,5~ч, сон 8,5~ч, период 24,0~ч.\tS
\item Постоянные времени веществ сна: подъём 18,2~ч, спад 4,2~ч. Спад втрое с~лишним быстрее подъёма — сон восстанавливает быстрее, чем бодрствование нагружает.\tS
\item В~исходной статье на~оба направления стояла одна постоянная 45~ч. Тогда пила получается почти прямой, а~главное — сон восстанавливает ровно с~той~же скоростью, с~какой бодрствование нагружает, чего быть не~может. После замены размах веществ сна вышел 0,527 против 0,50 у~авторов двухпроцессной модели — и~это совпадение никто не~подгонял.\tD
\item Разброс длительности сна от~ночи к~ночи у~людей: 1,29~ч по~актиграфии (26\,121 ночь), 1,12~ч по~ЭЭГ, 0,54~ч по~самоотчёту.\tP
\end{itemize}

\begin{ineq}
Пара «тормози друг друга» даёт S-образную кривую стационарных состояний; средняя ветвь неустойчива, крайние устойчивы. Переходы — седло-узловые бифуркации на~концах средней ветви. Медленная переменная (вещества сна) ползёт вдоль кривой и~сбрасывает систему со~складки; это релаксационный осциллятор, подстроенный внешним периодическим членом. Важно, что положение складок — свойство быстрой подсистемы, а~не~порог, который можно было~бы назначить.
\end{ineq}

% =====================================================================
\section{От одной клетки к популяции}\label{sec:pop}
\fig{fig_population.png}{Одна клетка имеет два состояния и~ничего между ними. У~клеток разная стойкость, поэтому доля повреждённых меняется плавно — но~петля гистерезиса сохраняется.}

До~сих пор речь шла об~одной клетке. Но~тормозных клеток в~коре много, и~ведёт себя популяция иначе, чем любая из~её клеток.

Одна клетка с~двумя состояниями — это выключатель: либо здорова, либо повреждена, промежуточного нет. Это видно и~в~расчёте: за~всё время счёта одиночная клетка проводила в~середине шкалы повреждения около одного процента времени. То~есть медленного ухудшения у~неё не~бывает вовсе: только скачок.

Если клетки чуть различаются по~стойкости — а~они различаются, — то~при растущей нагрузке они ломаются не~разом, а~по~очереди. Доля сломанных растёт плавно. При этом петля гистерезиса никуда не~девается: чтобы клетки вернулись, нагрузку надо опустить ниже, чем та, при которой они ломались. Получается плавная величина, сохранившая память, — то, чего у~одной клетки нет.

Здесь есть тонкость, из-за которой нельзя просто усреднить. Если взять популяцию независимых клеток и~посчитать среднее, плавность получится «из~ничего» — просто от~усреднения, без всякого механизма. Поэтому клетки в~модели связаны через то, через что они связаны в~коре: через активность сети. Меньше работающих тормозных клеток — слабее торможение — выше возбуждение — выше нагрузка на~оставшихся. Круг замыкается на~самом деле, а~не~на~словах.

\facts
\begin{itemize}
\item Среднее по~популяции двухустойчивых единиц с~разбросанными порогами меняется плавно и~сохраняет гистерезис.\tP
\item Собственная разнородность популяции не~сворачивается в~одно уравнение среднего поля: нужна доля сломанных, а~не~среднее состояние.\tP
\item Расчёт: у~одиночной клетки на~середину шкалы приходится 0,9–1,3~\% времени; у~популяции из~32 клеток — около 5~\%.\tD
\item Отрицательный результат. Плавность даёт именно разброс стойкости, а~не~связь через сеть: при самом узком разбросе она почти исчезает. Надежда, что круг «активность → нагрузка → повреждение → торможение → активность» сам создаст промежуточные состояния, не~подтвердилась.\tD
\end{itemize}

\begin{ineq}
Формально: популяция одинаковых бистабильных элементов имеет ту~же бифуркационную диаграмму, что и~один элемент. Разброс параметра размывает обе складки в~интервал, и~доля элементов в~верхнем состоянии становится функцией распределения порогов — то~есть ступенчатой функцией с~шагом $1/J$, монотонной и~гистерезисной. Общая связь через среднее поле добавляет к~этому положительную обратную связь, то~есть может сузить интервал обратно или даже вернуть разрывный переход.
\end{ineq}

% =====================================================================
\section{Кора в петле}\label{sec:loop}
Чтобы круг замкнулся, нужна модель коры — хотя~бы грубая. Берём две группы клеток: возбуждающие и~тормозные PV-клетки. Каждая описана одним числом — средней частотой разрядов. Возбуждающие подстёгивают сами себя и~тормозных, тормозные гасят и~тех и~других.

Повреждение входит сюда одной величиной: выходные веса тормозных клеток умножаются на~множитель, который падает с~долей повреждённых. При полной потере торможение опускается до~60~\% от~исходного — это число взято из~работы, где такую потерю моделировали.

Что даёт эта связь: когда торможение падает до~60~\%, активность возбуждающих клеток поднимается на~19~\%. На~столько~же поднимается и~нагрузка на~оставшиеся тормозные клетки. Это и~есть замкнутый круг, и~его сила — не~подобранное число, а~следствие двух взятых из~литературы.

\facts
\begin{itemize}
\item Модели нейронных масс, где тормозной вес — сама по~себе переменная со~своим уравнением, в~литературе есть; падение торможения до~60~\% от~исходного при потере PV-клеток — тоже.\tP
\item Множитель ставится на~все выходные веса тормозной группы сразу, а~не~на~один: так сделано в~работе, где разбирали именно потерю PV.\tP
\item Моделей, связывающих обмен веществ в~клетке с~нейронной массой, найти не~удалось. Круг «перевозбуждение → повреждение обмена → меньше торможения → перевозбуждение» в~литературе описан словами, без уравнений.\tS
\end{itemize}

\begin{note}
Кора здесь взята в~стационарном состоянии: её собственные времена — десятки миллисекунд, то~есть в~тысячи раз быстрее всего остального. Поэтому гамма-ритма в~этой версии нет намеренно: нужна средняя частота разрядов, от~неё зависит обмен. Режим с~колебаниями — отдельная задача.
\end{note}

% =====================================================================
\section{Закон сохранения разрывов}\label{sec:law}
\fig{fig_gaps.png}{Эпизод начинается, когда нагрузка выбрасывается выше верха окна, и~кончается, когда уходит ниже низа. Сумма двух скачков не~зависит от~того, где в~окне сидит средняя нагрузка.}

Это главный результат, который модель выдала не~как число, а~как запрет.

Картина болезни, которую надо получить: эпизоды редки и~длятся от~часов до~нескольких суток. Значит, сорваться должно быть трудно, а~вернуться — легко. В~модели и~то и~другое делает один и~тот~же случайный шум нагрузки: вверх перебрасывает через верхнюю границу окна, вниз возвращает под нижнюю.

И~тут выясняется простая вещь. Расстояние вверх и~расстояние вниз, если мерить их в~логарифме нагрузки, в~сумме дают ширину окна — и~ничего больше. Сдвинешь среднюю нагрузку вниз — до~срыва станет далеко, зато до~возврата близко; сдвинешь вверх — наоборот. Сумма не~меняется.

У~типичного набора параметров эта сумма составляет около девяти стандартных отклонений шума. А~картина болезни требует трёх отклонений вверх (чтобы срывы были редки) и~полутора вниз (чтобы возврат занимал сутки) — всего четырёх с~половиной. Девять в~четыре с~половиной не~помещаются.

Вывод жёсткий и~от~параметров не~зависит: \textbf{пока срыв и~возврат делает один и~тот~же шум, редких эпизодов с~быстрым возвратом не~будет.} Можно крутить любые числа — не~поможет.

Что из~этого следует. Нужна несимметрия: срыв и~возврат должны управляться разными вещами. Два пути видны уже сейчас, и~оба проверяемы.

\begin{itemize}
\item \textbf{Односторонний триггер.} Всплеск серотонина поднимает нагрузку и~никогда её не~опускает. Тогда разрыв вверх перекрывается им, а~разрыв вниз остаётся на~обычном шуме, и~два расстояния меряются разными распределениями.
\item \textbf{Переменная защита.} Если после нагрузки клетка усиливает ремонт, нижняя граница окна сама подтягивается вверх, и~возврат происходит без большого выброса. Тогда закон перестаёт связывать с~другого конца — он выведен в~предположении, что границы окна стоят на~месте.
\end{itemize}

\facts
\begin{itemize}
\item Активация рецептора 5-HT2A поднимает частоту разрядов возбуждающих клеток коры примерно вдвое; путь подтверждён тем, что эффект снимается и~блокадой рецептора, и~блокадой внутриклеточного звена. Это фармакологический максимум, а~не~обычный всплеск.\tP
\item Длительность острого действия веществ, работающих через этот рецептор: 6–7~ч и~10–11~ч, начало через 0,4–0,8~ч.\tP
\item Фермент ремонта пероксиредоксина наращивается активностью: мРНК в~13,5 раза за~4~часа, белок — к~суткам. Скорость обратного спада ни~в~одной из~прочитанных работ не~измерена.\tS
\end{itemize}

\subsection{Уточнение, которое не~подтвердилось}
В~версии~3.1 здесь стояло объяснение: будто эпизод не~кончается потому, что случайный провал нагрузки слишком короток — клетка чинится около суток, а~быстрая составляющая шума держится два часа. Из~него следовали два предсказания: у~зависимости длительности эпизода от~времени корреляции шума должен быть минимум около времени починки, и~удлинение быстрой составляющей должно эпизод укоротить.

Проверено отдельным прогоном. **Оба предсказания оказались неверны.** Минимума нет вовсе: длительность растёт монотонно, и~чем короче шум, тем короче эпизод вплоть до~края перебора. Удлинение быстрой составляющей эпизод не~укорачивает, а~удлиняет — с~3,4 до~7,9~суток. Объяснение снято.

Что осталось вместо него — измеренная зависимость, без объяснения:
\[\text{длительность} \;\propto\; \tau^{1/4}\,\sigma^{2/3},\]
где $\tau$ — время корреляции шума, $\sigma$ — его сила. Показатель при~$\tau$ вышел одинаковым для~обеих составляющих, медленной и~быстрой: 0,254 и~0,243 при~разбросе значений в~40 и~33 раза. Такое совпадение само просит объяснения, но~давать его до~проверки было~бы повторением той~же ошибки.

Практический вывод, который из~этого следует сразу. При~времени корреляции в~шесть часов эпизоды выходят по~1,5~суток, и~больше половины из~них короче суток — то~есть именно та~картина, которой мы добивались. Но~предвестник сна, описанный в~разделе~\ref{sec:supp}, требует шума с~временем в~несколько суток: иначе ему просто не~из~чего взяться. Одним источником случайности обе вещи получить нельзя. Значит, эпизод должен заканчиваться не~по~воле шума, а~сам — и~это ещё один довод в~пользу переменной защиты.

\begin{ineq}
Пусть $W$ — отношение границ окна, $q$ — положение средней нагрузки в~нём по~логарифмической шкале, $F$ — множитель обратной связи при полном повреждении, $P$ и~$T$ — суточные пик и~впадина. Тогда
\[\underbrace{\ln\frac{W^{1-q}}{P}}_{\text{до срыва}} + \underbrace{\ln\big(W^{q} F T\big)}_{\text{до возврата}} = \ln\frac{W F T}{P},\]
и~правая часть от~$q$ не~зависит. Утверждение элементарное — это сложение двух логарифмов, — но~оно ограничивает любую модель такого устройства: бистабильная единица под зашумлённой нагрузкой не~может одновременно иметь высокий барьер входа и~низкий барьер выхода. Отсюда и~требование несимметричного механизма.
\end{ineq}

% =====================================================================
\section{Что модель уже не смогла}\label{sec:fail}
Отрицательные результаты здесь перечислены отдельно и~намеренно. Каждый из~них — это проверка, записанная до~счёта, которая не~прошла; каждый сузил круг возможных объяснений сильнее, чем любой из~успехов.

\begin{enumerate}
\item \textbf{Суточный размах — не~ворота, а~ручка громкости.} Ожидалось, что есть пороговая величина суточного размаха нагрузки, ниже которой срывов почти нет, а~выше — есть. Получился плавный рост без перелома: при размахе от~0 до~40~\% доля повреждавшихся выросла всего в~1,5 раза. Размах меняет риск, но~ничего не~включает.\tD
\item \textbf{Плавность — от~разброса, а~не~от~сети.} Надежда была, что связь через активность сети сама создаст промежуточные состояния между «здорова» и~«повреждена». Нет: при узком разбросе стойкости клеток плавность почти исчезает (рис.~\ref{fig:res}а).\tD
\item \textbf{Предвестника сна без медленной переменной не~бывает.} Пока единственный источник случайности менялся за~2,4~часа, никакого изменения сна за~сутки до~срыва не~возникало — и~не~могло: в~модели не~было ни~одной величины, живущей на~масштабе дней. Предвестник появился только после того, как такая величина была добавлена.\tD
\item \textbf{Длительность эпизода осталась неверной.} В~модели эпизод длится около шести суток, тогда как по~наблюдениям — от~часов до~нескольких. Это прямое следствие закона из~раздела~\ref{sec:law}, и~исправляется оно не~параметрами.\tD
\item \textbf{Замена «доли сломанных» на~«среднее повреждение» не~изменила ничего.} Довод был сильный: доля равна нулю, пока не~сломается первая клетка, а~среднее растёт сразу. На~деле в~узком окне клетки ломаются почти одновременно, и~обе величины совпали с~точностью до~половины процента.\tD
\item \textbf{Серотониновый всплеск оказался не~спусковым крючком, а~той~же ручкой громкости.} Проверка была прямая: редкие односторонние всплески, поднимающие нагрузку, должны были перекрыть разрыв вверх и~оставить возврат обычному шуму. Из~четырёх частей правила не~прошла ни~одна. Срывов, случившихся вслед за~всплеском, оказалось 20~\% вместо ожидавшихся 80; длительность эпизода не~изменилась вовсе — 5,8~суток без всплесков и~5,2 с~ними. Всплески добавляют срывов (доля задетых наборов растёт с~0,59 до~0,73 при учащении всплесков), но~не~меняют ни~одного свойства самого эпизода.\tD
\item \textbf{Эпизод не~помнит, насколько глубоко зашёл.} Проверялось, живёт~ли эпизод своей жизнью после того, как нагрузка вернулась к~норме. Оказалось, что почти нет: длительность эпизода совпадает со~временем, которое нагрузка провела выше верхней границы окна, с~ранговой связью 0,94, а~связь восстановительного хвоста с~достигнутой глубиной повреждения — всего 0,11. Хвост при~этом занимает четверть эпизода и~длится около половины суток независимо от~того, сколько клеток сломалось.\tD
\item \textbf{Первое правило оказалось почти пустым.} Проверка на~опережение повреждения относительно защиты прошла, но~затем выяснилось, что звено первого порядка выполняет её почти всегда. Опережение на~треть периода, которое наблюдают в~опытах, этому ядру недоступно ни~при каких параметрах.\tD
\end{enumerate}

\begin{note}
Дважды за~работу правило проходило «вхолостую»: формально выполнялось, а~при разборе оказывалось, что сработал посторонний эффект. Один раз сравнивались разные наборы вместо одного и~того~же; другой раз контроль опирался на~единственный набор, и~статистика просто не~посчиталась. Оба случая вскрылись только потому, что к~каждому правилу заранее была записана не~только проверка, но~и~способ её обмануть.
\end{note}

% =====================================================================
\section{Дополнительные результаты}\label{sec:supp}
\figp{fig_results.png}{\label{fig:res}Четыре результата, которые видно без подготовки. Числа — медианы по~нескольким сотням случайных наборов параметров.}

Устройство работы такое: прежде чем что-либо считать, формулируется правило — что именно должно получиться и~с~каким числом, — и~только потом запускается счёт. Правило нельзя переписать после того, как результат увиден; можно лишь признать его негодным и~объяснить, почему. Ниже — все правила проекта и~что из~них вышло.

{\small
\begin{tabular}{@{}R{0.05}R{0.40}R{0.31}R{0.17}@{}}
\toprule
 & \textbf{что проверяли} & \textbf{что вышло} & \textbf{вывод}\\
\midrule
Ф1 & повреждение опережает защиту на~1/8–1/4 периода & 0,162 & выполнено, но~правило признано почти пустым\\
Ф2а & срывы копятся ко~второй половине бодрствования & 0,88 при~случайном 0,50; аналитическое ожидание 0,82 & выполнено, и~совпало с~выкладкой\\
Ф2б & есть порог по~суточному размаху & рост в~1,5 раза, перелома нет & \textbf{не~выполнено}\\
Ф2б$'$ & кривая ложится на~гладкое предсказание без~ворот & худшее отклонение 0,55 при~допуске 0,04 & \textbf{не~выполнено}: ни~ворот, ни~гладкой опасности\\
Ф2в & сон редко обрывает эпизод & 0,51 при~равномерном ожидании 0,35 & \textbf{не~выполнено}: сон обрывает \emph{чаще} случайного\\
Ф3а & повреждение мешает заснуть и~удлиняет эпизод & в~1,67 раза при~пороге 1,5 & выполнено\\
Ф3а$'$ & ...и~ровно настолько, насколько удлиняется бодрствование & наблюдение 1,67 против 1,10 у~переключателя & непроверяемо; усиление есть и~требует объяснения\\
Ф3в & эпизоды короче суток редки & 0,13 по~всем эпизодам при~пороге 0,50 & выполнено\\
Ф4а & сон портится ещё до~порога срыва & 1,085 при~пороге 1,10 и~ожидании переключателя 1,114 & \textbf{не~выполнено}, и~ниже механического ожидания\\
Ф4б$'$ & предвестник есть по~собственному уровню набора & 1,055 при~пороге 1,05 & выполнено\\
Ф5 & доля повреждённых стала плавной & 0,038 при~пороге 0,05 и~вилке выкладки 0,069–0,139 & \textbf{не~выполнено, и~ниже вилки}: гистерезис убирает промежуточные состояния\\
Ф6а & бодрствование растёт перед срывом & 1,055, в~механической вилке & выполнено\\
Ф6б,г & наклон кривой предвестника & $-0{,}62$ по~восьми ячейкам, MDE 0,94 & непроверяемо при~любой выборке: дробить ячейки\\
Ф6в$'$ & медленный шум поднимает предвестник (парно) & $p=0{,}97$ & \textbf{не~выполнено}: шум ни~при~чём\\
Ф7а & эпизод не~длиннее двух суток & 4,39~суток & \textbf{не~выполнено}\\
Ф7а$'$ & длительность идёт по~закону Ф16 вне подгонки & промах 0,22~суток при~допуске 1,5 & выполнено\\
Ф7б & есть эпизоды короче суток & 0,28 по~всем эпизодам при~пороге 0,30 & \textbf{не~выполнено}\\
Ф8а & срыв вызывается всплеском серотонина & 0,24 при~пороге 0,80 и~ожидании выкладки 0,36 & \textbf{не~выполнено}: порог был требованием к~отбору\\
Ф8б,в & всплеск меняет свойства эпизода & 7,2~суток; коротких 0,12 & \textbf{не~выполнено}\\
Ф10а & хвост восстановления связан с~глубиной повреждения & переписано под~нынешнюю модель, ждёт прогона & записано до~счёта\\
Ф10б,в & эпизод~— это время высокой нагрузки, а~не~ремонт & предварительно 0,87 и~хвост 7~\% & записано до~счёта\\
Ф11 & индукция закрывает эпизод & связь силы индукции с~длительностью слабее порога & \textbf{не~выполнено}\\
Ф12 & защита создаёт глухой период & по~промежуткам между эпизодами~— нет & \textbf{не~выполнено}\\
Ф13 & ответ не~зависит от~числа клеток & размах 0,21 и~0,24 при~потолке шума 0,055 и~0,042 & \textbf{не~выполнено}: $J$~— параметр, а~не~разрешение\\
Ф14 & память защиты маскируется шумом & две части из~трёх & частично\\
Ф15 & глухой период есть и~он не~от~защиты & $t_r=1{,}0\dots1{,}4$~суток, оба контроля прошли & выполнено\\
Ф16 & закон длительности предсказывает вне подгонки & промах 10~\% при~допуске 2~SE & выполнено\\
Ф16в & закон аффинный, а~не~насыщающийся & прямая обошла обе насыщающиеся формы & выполнено\\
Ф18б$'$ & предвестник делает память защиты (парно) & падает у~107~наборов, $p\le0{,}05$ & выполнено\\
Ф18б$''$ & ...и~делает его целиком & уходит лишь пятая часть превышения & \textbf{не~выполнено}: есть третья медленная величина\\
Ф19 & прайминг доводит эпизод до~цели & 1,46~суток, граница 1,665, $p=5{,}2\cdot10^{-6}$ & \textbf{выполнено, все четыре части}\\
Ф20 & ритм ломается при~сильной связи & лестница захвата 1:1 → 5:4 → 3:2 → 2:1 → 3:1 → 4:1 → сна нет & выполнено; это лестница, а~не~каскад удвоений\\
Ф20д & есть режимы с~чередованием длин циклов & узоры 25,\,24,\,24,\,47 и~47,\,25 и~46,\,26,\,24~ч & выполнено; рабочая область до~них не~достаёт\\
Ф21 & третья медленная величина~— заполнение популяции & записано до~счёта & ждёт прогона\\
Ф22 & $J$ входит через эффективный разброс сетки & арифметика сетки подтверждена; сверка подменой сигмы записана & ждёт прогона\\
\bottomrule
\end{tabular}}

\subsection{Сверка с~опубликованными числами}
Отдельно от~правил идут проверки, где модель обязана воспроизвести чужой результат. Это не~успех модели, а~условие допуска к~работе: если не~сошлось, считать дальше нельзя.

{\small
\begin{tabular}{@{}R{0.37}R{0.28}R{0.30}@{}}
\toprule
\textbf{величина} & \textbf{наш расчёт} & \textbf{в~статье}\\
\midrule
складки переключателя сна & 1,450 и~2,463~мВ & 1,45 и~2,46~мВ\\
бодрствование / сон / период & 15,50 / 8,50 / 24,00~ч & 15,5 / 8,5 / 24,0~ч\\
размах давления сна & 0,527 & 0,50\\
разброс длительности сна & 0,67~ч & 1,12–1,29~ч (с~режимом дня)\\
шаг счёта против точного решателя & расхождение 3\,$\cdot$\,10$^{-7}$ & —\\
\bottomrule
\end{tabular}}

Размах давления сна в~третьей строке не~подгонялся вовсе: постоянные времени взяты из~одной работы, устройство переключателя — из~другой, а~совпадение получилось само. Это и~есть самая весомая из~проверок.

% =====================================================================
\section{Медленная защита: что добавилось и что она сделала}\label{sec:slow}
\figp[0.88\textwidth]{fig_slow.png}{\label{fig:slow}Четыре результата версии~4. \textbf{а}~— две ступени защиты: сигнал индукции быстрый, белок медленный, и~половина уровня набирается только к~третьим суткам. \textbf{б}~— закон длительности эпизода: он оказался прямой, а~не~насыщающейся кривой, и~проверочная точка вне подгонки это подтвердила. \textbf{в}~— лестница захвата: число вращения стоит на~рациональных значениях и~перескакивает; во~врезке~— узор чередования длин циклов на~дробной ступени 5:4. \textbf{г}~— число клеток входит в~модель как эффективный разброс стойкости.}

Закон сохранения разрывов (раздел~\ref{sec:law}) оставил ровно два выхода, и~второй из~них~— переменная защита: если клетка после нагрузки усиливает ремонт, нижняя граница окна подтягивается вверх сама, и~возврат не~требует большого выброса шума. Этот выход мы и~достраивали.

\textbf{Что именно добавлено.} Фермент ремонта пероксиредоксина (сульфиредоксин) наращивается не~мгновенно: сначала накапливается сигнал индукции, потом по~нему синтезируется белок. Это две ступени первого порядка подряд:
\[\dot m=\frac{1}{\tau_m}\Big(\beta\frac{u}{K_B+u}-m\Big),\qquad \dot s=\frac{1}{\tau_s}\,(m-s),\]
где $u$~— перекись, $m$~— сигнал, $s$~— наработанный белок, который прибавляется к~постоянной защите. Стационар тот~же, что был раньше: окно, его границы и~весь закон сохранения разрывов не~меняются. Меняется только ход во~времени~— и~именно он решает всё.

\textbf{Откуда постоянные.} Время первой ступени взято по~мРНК: 8,1–9,6~ч в~трёх независимых обработках.\tS\ Время второй~— по~белку из~сырых данных той~же лаборатории: пик на~18-м часу, а~постоянная спада в~опубликованных работах не~измерена ни~разу, и~по~разным способам обработки выходит от~24 до~88~ч.\tS\ Рабочее значение 63~ч~— середина этого разброса, и~это наше допущение, а~не~измерение.\tD\ Сила индукции разыграна по~наборам: измеренная кратность наработки белка идёт $\times$2,1 / 2,7 / 4,2 / 7,0 при~четырёх дозах окислителя.\tS

\textbf{Что из~этого вышло~— и~это главный отрицательный результат версии.} Длительность эпизода оказалась \emph{прямой} функцией памяти защиты (рис.~\ref{fig:slow}б):
\[\text{длительность}\;=\;1{,}28\;+\;0{,}97\,\tau_{\rm б}\quad\text{(суток)},\]
без всякого насыщения. Форма выбиралась не~на~глаз: три кандидата подгонялись по~точкам с~памятью ниже 100~ч, а~проверялось предсказание в~точке 150~ч, в~подгонку не~входившей. Прямая выиграла у~обеих насыщающихся и~по~невязке подгонки, и~по~этой проверке.

Отсюда следствие, которое закрывает целую дорогу. Чтобы эпизод уложился в~клинические двое суток, нужна память защиты около 12~часов~— то~есть \emph{меньше, чем у~первой ступени}, и~в~несколько раз меньше любого из~измеренных значений (24–88~ч). Пассивная защита не~промахивается мимо цели на~волосок: она исключена в~2–10~раз. Усиление ремонта само по~себе эпизод не~укорачивает, а~удлиняет~— потому что чем дольше держится наработанный белок, тем дольше клетка сидит в~промежуточном состоянии, вместо того чтобы быстро свалиться обратно.

\begin{ineq}
Две ступени первого порядка подряд дают отклик на~ступень нагрузки
\[s(t)=1-\frac{\tau_s e^{-t/\tau_s}-\tau_m e^{-t/\tau_m}}{\tau_s-\tau_m},\]
то~есть всегда с~перегибом: мгновенного отклика нет ни~при каких постоянных. При~$\tau_m=12$~ч и~$\tau_s=63$~ч половина уровня набирается к~57~часам (рис.~\ref{fig:slow}а). Это и~есть причина, по~которой ускорять надо не~первую ступень, а~\emph{усиление}: производная отклика по~$\tau_m$ в~рабочей точке мала (на~одной клетке ускорение дало $-1\dots-5$~\%), а~по~амплитуде~— велика ($-52$~\%). Предсказание было записано до~счёта и~подтвердилось.
\end{ineq}

\begin{note}
Из~того, что защита в~модели эпизод удлиняет, \emph{не} следует, что сульфиредоксин вреден. Следует одно: защита, которая только чинит повреждённое, картину болезни не~объясняет. Объяснить её может лишь защита, которая срабатывает \emph{раньше и~сильнее}~— то~есть прайминг (раздел~\ref{sec:prime}).
\end{note}

% =====================================================================
\section{Прайминг: единственная дорога, которая довела до~цели}\label{sec:prime}
Прайминг~— это память о~прошлом воздействии: клетка, уже встречавшая окислитель, отвечает на~следующее воздействие сильнее и~быстрее. В~модель оно вставлено одной строкой: тем~же сигналом, что движет индукцию, набирается след $p$, и~сила индукции умножается на~$(1+A\,p)$. Потолок $A$ взят по~измерению~— втрое.

\facts
\begin{itemize}
\item При~повторном воздействии через 24~часа отклик регулятора ниже, а~наработка самого фермента \emph{выше}~— до~трёх раз, и~время до~полумаксимума падает примерно с~8 до~4~часов. Та~же лаборатория и~та~же система, что дала нам сырые данные по~кинетике.\tS
\item Эффект зависит от~главного регулятора окислительного ответа и~не~зависит от~его вспомогательного партнёра: то~есть это именно усиление сигнала, а~не~побочный эффект.\tS
\item Независимо, в~другой ткани: накопление защитного фермента при~повторных воздействиях прослежено на~четырёх сутках и~избирательно по~генам.\tP
\end{itemize}

\textbf{Что получилось.} Впервые за~проект длительность эпизода вошла в~клиническую цель: медиана 1,46~суток при~верхней границе 1,665 (односторонняя оценка по~порядковой статистике), доля коротких эпизодов 0,363, парный критерий против контроля без прайминга $p=5{,}2\cdot10^{-6}$. Все четыре части правила, записанного до~счёта, выполнены~— и~повторно выполнены на~выборке в~2000~наборов, вшестеро большей.

Важно, чем это \emph{не} является. Прайминг не~подгонялось: его амплитуда взята потолком из~измерения, его постоянная времени перебиралась, а~ускорение первой ступени~— вторая половина измеренного эффекта~— было проверено отдельно и~не~дало почти ничего, как и~предсказывалось до~счёта. Дорога осталась одна, и~она прошла.

% =====================================================================
\section{Предвестник: кто его делает}\label{sec:prodrome}
За~несколько суток до~срыва бодрствование удлиняется~— в~модели это видно как подъём кривой «длина бодрствования против срока до~срыва», нормированной на~собственный здоровый уровень каждого набора. Вопрос, кто его делает, решался парным сравнением: одни и~те~же наборы, одно сид, разница только в~одном звене.

\begin{itemize}
\item \textbf{Медленный шум нагрузки~— не~он.} При~его включении высота предвестника не~растёт (парный знаковый критерий: подъём у~131 набора из~выборки, $p=0{,}97$). Шум вводился под~обоснование «в~модели нет ни~одной величины, живущей на~масштабе дней»; это перестало быть правдой, как только появилась двухступенчатая защита. Шум выведен из~рабочего набора.
\item \textbf{Память защиты~— частично он.} При~её отнятии высота падает у~107~наборов из~проверенных, $p\le0{,}05$. Но~уходит только \emph{пятая часть} превышения.
\end{itemize}

Значит есть третья медленная величина. Поиск сузился до~одного кандидата простым перебором всех времён модели: возврат перекиси и~релаксация нагрузки дают 0,25–0,5 и~0,03–0,05~суток; возврат $x$~— порядка 0,01~суток (и~это подтверждено прямым измерением: у~73~\% наборов $x$ возвращается за~сутки); первая ступень защиты~— 0,5~суток; вторая~— 2,6~суток и~уже исключена как единственная причина. \emph{Других времён в~уравнениях нет.}

Отсюда гипотеза, записанная как правило Ф21: третья медленная величина~— не~постоянная времени, а~\textbf{накопление}. Популяция из~$J$ клеток заполняется повреждёнными по~одной, и~характерное время заполнения задаётся не~константой, а~\emph{редкостью} перехода, то~есть расстоянием нагрузки до~барьера в~сигмах шума. У~такой величины нет своего~$\tau$~— поэтому перебор постоянных её найти и~не~мог.

% =====================================================================
\section{Из~чего состоит эпизод}\label{sec:tail}
Все прочие правила меряют длительность \emph{между} наборами. Но~для~выбора механизма важнее, из~чего длительность складывается \emph{внутри} набора. Эпизод делится на~две части: время, пока нагрузка держится выше верха окна, и~\textbf{хвост}~— время от~последнего такого момента до~конца эпизода, когда нагрузка уже нормальная, а~клетки ещё повреждены.

Это различие решает, за~что вообще может зацепиться защита. Если хвост велик, лимитирует ремонт, и~ускорять надо его. Если хвоста почти нет, эпизод длится ровно столько, сколько держится высокая нагрузка, и~ремонт не~лимитирует вовсе.

На~прежней версии модели (ещё без двухступенчатой защиты) хвост занимал четверть эпизода, а~связь длительности со~временем высокой нагрузки была 0,94. Предварительная проверка на~\emph{нынешней} модели даёт картину резче: связь 0,87, хвост~— около 7~\% эпизода.\tD\ Если это подтвердится на~полной выборке, вывод будет жёстким: \textbf{прайминг укорачивает эпизод не~тем, что чинит быстрее, а~тем, что поднимает верх окна} и~тем сокращает само время высокой нагрузки. Это проверяемое предсказание, и~оно записано до~счёта.

\begin{note}
Прежнее измерение хвоста сделано на~модели без двухступенчатой защиты, и~переносить его на~нынешнюю нельзя: именно защита двигает границы окна. Поэтому правило переписано, а~регистратор эпизодов \emph{сверяется побитово} с~основным счётом на~тех~же строках и~том~же сиде: если копия шага разойдётся хоть в~одном эпизоде, правило объявляет себя непроверяемым и~чисел не~печатает.
\end{note}

% =====================================================================
\section{Сколько клеток в~популяции~— и~почему это оказался параметр}\label{sec:J}
Число клеток $J$ до~недавнего времени выбиралось по~удобству счёта: 32 в~одних прогонах, 64 в~других. Предполагалось, что это просто разрешение. Оказалось, что это параметр, и~входил он тремя разными путями.

\begin{enumerate}
\item \textbf{Хвост распределения стойкости.} Множители брались равновероятными квантилями $p_j=(j+1/2)/J$, поэтому крайний квантиль уползает с~ростом $J$: чем больше клеток, тем уязвимее самая уязвимая. \emph{Закрыто}: сетка вероятностей закреплена по~краям, и~при~$J=32$ она совпадает с~прежней до~последнего знака, так что старые прогоны остаются в~силе.
\item \textbf{Квантование порога эпизода.} Эпизод определялся как «доля повреждённых выше $1/2$», а~доля принимает только значения $j/J$, поэтому настоящий порог равен $0{,}625$ при~$J=8$ и~$0{,}516$ при~$J=64$. \emph{Обойдено}: вердикты выносятся по~непрерывной переменной, у~которой порог не~квантован.
\item \textbf{Реализованный разброс.} А~вот это нашлось только на~выборке в~2000~наборов, когда шум измерений упал ниже эффекта. Края сетки закреплены и~действительно почти не~ползут (наименьший множитель 0,8002 при~$J=8$ и~0,8029 при~$J=512$), но~\emph{дисперсия} конечной сетки меньше заданной и~подходит к~ней только в~пределе: $\mathrm{std}(\ln v)=0{,}1226\,/\,0{,}1060\,/\,0{,}0980\,/\,0{,}0943$ при~$J=8/16/32/64$ и~$0{,}0910$ в~пределе (рис.~\ref{fig:slow}г). От~$J=16$ к~$J=64$ это 11~\% по~эффективной сигме разброса, а~модель к~ней крайне чувствительна.
\end{enumerate}

Третий путь объясняет найденную зависимость количественно и~указывает поправку: растянуть сетку так, чтобы реализованный разброс равнялся заданному при~любом~$J$. Решающая проверка записана до~счёта: блок при~$J=32$ с~подменённой сигмой обязан воспроизвести блок при~$J=16$ с~обычной. Если воспроизведёт~— путь объяснён целиком.

\begin{ineq}
Общий урок, не~про~нашу модель. Дискретизация непрерывного распределения $J$~точками может \emph{точно} сохранять среднее (нормировкой) и~\emph{точно} закреплять края (выбором крайних квантилей) и~при~этом не~сохранять дисперсию: для~равномерной по~вероятности сетки $\mathrm{Var}\to\sigma^2$ лишь как $O(1/J)$ сверху. Если модель чувствительна к~дисперсии разброса~— а~здесь чувствительность по~логарифму около~2,5,~— то~число точек сетки становится параметром модели, и~сравнивать прогоны с~разным $J$ нельзя: это сравнение двух разных моделей.
\end{ineq}

% =====================================================================
\section{Как ломается ритм}\label{sec:bifur}
Отдельный вопрос, поставленный независимо от~редокса: что происходит с~чередованием сна и~бодрствования, когда повреждение всё сильнее тормозит сонное ядро. Ожидание было такое: при~приближении к~критическому значению должен появляться хаос~— то~есть циркадный ритм ломается \emph{при~сохранённом} 24-часовом внешнем задатчике. Считается это дёшево: переключатель детерминированный, клеток и~редокса не~нужно вовсе.

Правильная наблюдаемая здесь~— \textbf{число вращения}: сколько суток внешней силы приходится на~один цикл сна и~бодрствования. Ритм не~ломается плавно: число вращения стоит на~рациональных значениях в~конечных интервалах параметра и~перескакивает между ними (рис.~\ref{fig:slow}в):
\[1\!:\!1\;\to\;5\!:\!4\;\to\;3\!:\!2\;\to\;2\!:\!1\;\to\;3\!:\!1\;\to\;4\!:\!1\;\to\;\dots\;\to\;9\!:\!1\;\to\;\text{сна нет}.\]
Переходы приходятся на~$kD\approx4{,}77$; 4,78; 4,79; 5,09; 5,16 и~далее всё теснее: от~5,172 до~5,178~мВ умещаются сразу пять ступеней (5:1, 6:1, 7:1, 8:1, 9:1), а~потом сон исчезает совсем.

Отдельный прогон с~мелким перебором и~измельчением границ (рис.~\ref{fig:stairs}) показывает, что ступеней гораздо больше: \textbf{52 на~участке в~0,9~мВ, и~у~34 из~них длины циклов чередуются}. Между целыми захватами стоит плотное скопление дробных~— 5:4, 6:5, 4:3, 13:10, 11:8, 9:7, 7:5, 10:7, 16:11, 8:5, 12:7,~— а~между ними точки, где узор не~повторяется вовсе. Ширины ступеней падают с~ростом знаменателя на~много порядков: у~2:1 это 0,3~мВ, у~3:2 около 8$\cdot$10$^{-3}$, у~11:8 уже 5$\cdot$10$^{-6}$. Сплошным перебором такое не~поймать: приходится измельчать только границы.

\figp{fig_stairs.png}{\label{fig:stairs}Чёртова лестница целиком, по~отдельному прогону с~измельчением границ (около 60 циклов на~точку после сброса переходного процесса). \textbf{Сверху}~— весь участок: ровные ступени целых захватов 1:1, 2:1, 3:1 и~дальше до~9:1 вплотную перед исчезновением сна. \textbf{Снизу}~— крупный план той области, где ступеней больше всего: оранжевым показаны захваты с~\emph{чередованием} длин циклов, серыми точками~— режимы без повторяющегося узора. Во~врезке~— сам узор захвата 3:2: длинный цикл 47~ч, короткий 25~ч, и~так по~кругу.}

\textbf{Это лестница захвата, а~не~каскад удвоений.} В~каскаде удвоений период умножается на~два и~других значений не~принимает, а~постоянная Фейгенбаума описывает, как сходятся ступени. Здесь период прибавляет по~одной пропущенной ночи, и~между целыми ступенями стоят дробные~— это языки Арнольда, обычная картина захвата автоколебания внешней силой. Ссылаться на~постоянную Фейгенбаума здесь было~бы ошибкой.

\textbf{Самое интересное~— дробные ступени.} В~них длины циклов не~одинаковы: они идут строгим узором и~повторяются. При~$kD=4{,}772$~мВ узор из~четырёх циклов на~пять суток: 25, 24, 24, 47~ч~— трое почти обычных суток, затем одни удвоенные, и~снова (врезка на~рис.~\ref{fig:slow}в). При~4,779~— узор из~двух: 47 и~25~ч. При~4,783~— из~трёх: 46, 26, 24~ч. Проверено на~устойчивость: узор тот~же при~120 и~при~400~сутках счёта, то~есть это не~переходный процесс.

Это важно вот почему. Такое поведение~— не~«ритма нет», а~\emph{другой ритм}, с~собственным периодом в~несколько суток, и~именно так выглядит часть клинических картин нарушенного сна. Модель даёт его сама, без всякого шума и~без единого добавленного уравнения,~— просто как следствие захвата. Предсказание, записанное до~счёта: при~добавлении шума эти окна должны не~исчезнуть, а~размыться в~\emph{перемежающееся} чередование~— узор то~держится, то~срывается.

\textbf{Главное ограничение честное и~неприятное~— но~не~то, которое стояло здесь в~версии~4.} Потолок задаёт не~повреждение, а~сама связь: тождественно $kD\le k$, поэтому при~$k=2$~мВ ступень на~4,77~мВ недостижима \emph{при любом} повреждении. Нужна связь не~ниже 4,77~мВ, то~есть в~2,4~раза сильнее рабочей. При~$k=5$~мВ модель до~лестницы \emph{доходит}: у~наборов, свалившихся в~необратимое состояние, измерено $kD=4{,}93$ при~доле повреждённых 0,999 и~бодрствовании 0,974. Значит ограничение не~в~механизме, а~в~калибровке одной константы.

\textbf{Прежняя редакция называла два выхода; второй вычёркивается.} «Хроническое повреждение, которое держит долю сломанных клеток высокой», не~работает, и~это видно без~всякого счёта: у~сломанных наборов доля уже на~потолке, а~потолок равен~$k$. Остаётся один выход~— сила связи «повреждение~— сон». Прежние числа в~этом абзаце («вдвое», «вчетверо») были неверны и~не~сходились даже между собой; их корень разобран в~разделе~\ref{sec:prec} как седьмая статистическая ловушка.

\begin{note}
Проверка, записанная до~счёта, спрашивала не~то: она искала среди периодов значения 1, 2 и~4~— и~нашла их, но~нашла~бы и~в~лестнице захвата. Вердикт оставлен как есть, а~рядом записаны две части, которых не~хватало: есть~ли захват на~периоде, не~равном степени двойки, и~есть~ли режимы с~узором длиннее одного цикла. Отдельно пришлось исправить сетку перебора: ступени 5:4 и~3:2 имеют ширину около 0,005~мВ и~на~грубой сетке не~видны вовсе~— первый прогон увидел 1, 2, 4 и~не~увидел ни~одной дробной.
\end{note}

\begin{ineq}
Почему лестница, а~не~каскад, здесь и~должна была получиться. Переключатель~— релаксационный автоколебатель с~собственным периодом около 25~ч, суточный член~— внешняя сила периода 24~ч. Для~такой пары отношение периодов есть число вращения, и~оно захватывается на~рациональных $n\!:\!m$ в~интервалах конечной ширины (языки Арнольда), а~между ними меняется по~чёртовой лестнице. Ширина языка быстро падает с~ростом знаменателя~$m$: у~2:1 она около 0,3~мВ, у~3:2 и~5:4~— около 0,005, то~есть в~шестьдесят раз уже. Отличить лестницу от~каскада удвоений можно ровно одним признаком: наличием захвата на~периоде, не~равном степени двойки.
\end{ineq}

% =====================================================================
\section{Асимметричная защита}\label{sec:asym}
\figp[0.88\textwidth]{fig_v5.png}{\label{fig:v5}Четыре результата версии~5. \textbf{а}~— асимметричная защита: измеренный уровень ложится на~выкладку $w(r)$, порядок совпал точно. Размах привязан по~двум крайним точкам (в~кружках), три средние предсказаны. \textbf{б}~— рваность нагрузки: при~закреплённых среднем и~дисперсии частота срывов падает на~40~\%, а~доля повреждавшихся насыщается и~не~видит ничего; выбор измеряемой величины решил исход правила. \textbf{в}~— ширина полосы спасённых против доли повреждённых: два перебора дают противоположные знаки, то~есть полоса не~есть функция~$D$. \textbf{г}~— стоимость строки у~счёта по~одной операции удваивается, когда рабочий набор перестаёт влезать в~кэш, а~у~слитого ядра не~меняется вовсе.}

Вторая ступень защиты всё время была записана с~одной постоянной: $\dot s=(m-s)/\tau_s$, 63~часа и~на~подъём, и~на~спад. Биологически это два разных процесса~— индукция синтеза и~деградация белка,~— и~их равенство было допущением, которое мы просто ни~разу не~проверяли.\tD\ Версия~5 разводит их: $\tau_{\uparrow}$ на~подъёме, $\tau_{\downarrow}$ на~спаде, отношение $r=\tau_{\downarrow}/\tau_{\uparrow}$.

\textbf{Выкладка, дающая предсказание до~всякого счёта.} Вход $m$ при~суточном цикле с~хорошей точностью двухуровневый: высокий уровень держится долю $f$ суток, низкий~— остальное. Для утекающего интегратора с~разными постоянными, когда обе много больше суток, стационарное среднее стоит между уровнями в~точке
\[w(r)=\frac{f\,r}{f\,r+1-f},\]
то~есть $s_{\rm ср}=m_{\rm н}+w(r)\,(m_{\rm в}-m_{\rm н})$. Асимметрия превращает интегратор в~\emph{детектор пика} при~$r>1$ и~в~\emph{детектор дна} при~$r<1$. Доля~$f$ берётся из~прогона (она равна измеренной доле бодрствования, 0,705), а~не~назначается, поэтому в~предсказании нет ни~одного подогнанного числа.

\textbf{Что вышло.} При~$r=0{,}33/0{,}5/1/2/3$ выкладка даёт $w=0{,}441/0{,}544/0{,}705/0{,}827/0{,}878$, а~измеренный уровень защиты~— $0{,}0575/0{,}0644/0{,}0792/0{,}0958/0{,}1065$. Порядок совпал \emph{точно}: ранговая связь $+1$ при~точном перестановочном $p=0{,}0083$ (рис.~\ref{fig:v5}а). Следствие для срывов тоже подтвердилось: доля повреждавшихся падает с~0,486 до~0,457, критерий тренда даёт $z=1{,}67$, $p=0{,}048$, а~длительность эпизода укорачивается с~7,74 до~6,88~суток.

\textbf{Вывод, который меняет статус прежней модели.} Равенство постоянных подъёма и~спада было не~нейтральным упрощением, а~\emph{выбором режима}. Разные $r$ дают разную среднюю защиту при~одних и~тех~же мгновенных уравнениях, и~направление эффекта предсказуемо из~двух строк алгебры. Значит либо равенство надо обосновывать измерением, либо $r$ обязан остаться параметром модели наравне с~$\tau_s$.

\begin{note}
Количественная часть правила осталась непроверяемой, и~причина в~\emph{конструкции самого правила}, а~не~в~нехватке наборов. Допуск равносильности $3\,\mathrm{SE}=0{,}0241$ и~половина размаха кривой $0{,}0245$ почти совпали, то~есть зазор между порогом и~альтернативой 0,0004 при~MDE~0,02. Наборами это не~лечится: понадобилось~бы в~29~раз больше. Лечится плечом~— расширением перебора $r$ от~0,1 до~10, отчего размах кривой растёт в~1,7~раза. Это и~есть цена, записанная честно: мы знаем, \emph{чего} не~проверили и~\emph{сколько} это стоит.
\end{note}

% =====================================================================
\section{Три фазы: полоса спасённых}\label{sec:band}
Клетки одного набора отличаются ровно \emph{одним} скаляром~— множителем нагрузки~$v$; всё остальное у~них общее. Значит состояние клетки есть однопараметрическое семейство $x(t;v)$ при~общей истории нагрузки. Если~бы это семейство было монотонно по~$v$, повреждённые клетки были~бы всегда ровным хвостом $\{v\ge v^*\}$, доля повреждённых считалась~бы в~замкнутом виде $1-F_v(v^*)$, и~поклеточный блок был~бы не~нужен вовсе.

\textbf{Монотонности нет, и~ломает её медленная защита.} На~мелкой сетке в~600~узлов по~$v$ при~трёх профилях нагрузки, включая намеренно враждебный, измерено:
\begin{itemize}[leftmargin=1.1em,itemsep=1pt,topsep=2pt]
\item при~\emph{замороженной} второй ступени монотонность точная: ноль нарушений, наибольшее ровно $0{,}000$;
\item при~живой защите порядок рвётся на~11,9~\% шагов, наибольшее нарушение 0,973~— то~есть полный переворот: более уязвимая клетка цела, менее уязвимая рухнула;
\item прайминг не~причина, а~усилитель: при~$A=0$ нарушение остаётся, но~падает с~0,973 до~0,0184;
\item порядок не~сохраняется ни~в~одной координате: $s$~— 27,8~\% шагов, $m$~— 32,5~\%, след прайминга~— 32,7~\%, полная защита~— 100~\%.
\end{itemize}

\textbf{Механизм.} Более уязвимая клетка раньше окисляется, сильнее праймируется, копит больше защитного белка~— и~к~приходу большого выброса у~неё \emph{больше} защиты, чем у~менее уязвимой соседки. Это гормезис на~уровне популяции: часть клеток выживает именно потому, что раньше пострадала.

\textbf{Фаз оказалось три.} Пограничный слой не~размазан, а~связен: в~пробе его ширина держалась на~0,085 популяции как структура. В~боевом счёте при~$J=32$ ширина полосы (доля целых клеток, лежащих выше первой повреждённой) равна $0{,}00732\pm0{,}00042$~— семнадцать сигм от~нуля. Причина названа верно и~проверена: при~\emph{замороженной} защите полоса обращается в~\emph{ровный нуль}, при~включённом прайминге она в~1,82~раза шире ($0{,}01331$), а~короткая память второй ступени не~меняет её вовсе ($0{,}00712$ против $0{,}00732$).

\begin{ineq}
Форма полосы была предсказана геометрией: пограничный слой обязан следовать $\varphi(\Phi^{-1}(1-D))$, то~есть быть шире всего при~$D=1/2$. \textbf{Предсказание опровергнуто, причём дважды и~с~разными знаками.} По~перебору связи со~сном ранговая связь ровно~$-1$: полоса \emph{сужается} с~ростом доли повреждённых, с~$0{,}0135$ при~$D=0{,}10$ до~$0{,}0021$ при~$D=0{,}34$. По~независимому перебору суточного размаха, который для~этого вопроса не~открывался,~— ровно~$+1$: полоса \emph{расширяется}, с~$0{,}0095$ при~$D=0{,}064$ до~$0{,}0191$ при~$D=0{,}162$. Диапазоны перекрываются, и~при~одной и~той~же доле повреждённых ширина различается в~два раза в~зависимости от~того, \emph{какой ручкой} эту долю меняли (рис.~\ref{fig:v5}в). Утверждение о~форме снято целиком, как и~было записано заранее.
\end{ineq}

\textbf{Что это значит для~теории.} Замкнутая формула доли повреждённых через одну границу незаконна при~живой защите. Упрощать поклеточный блок можно только с~\emph{явным} пограничным слоем~— тремя конфигурациями, а~не~двумя. И~ширина этого слоя не~определяется наблюдаемой долей повреждённых: её задаёт что-то третье, скорее всего время, проведённое популяцией у~самого опрокидывания.

% =====================================================================
\section{Среднее поле: что из~него следует и~чего из~него не~выйдет}\label{sec:mf}
Структура модели выглядит как приглашение к~среднему полю: связь «все со~всеми» через два скаляра (доля повреждённых и~среднее повреждение), входящих только в~переключатель сна. Разбор по~методам даёт такую картину.

\textbf{Хартри, обычное среднее поле.} Не~работает, и~теперь известно почему: доля повреждённых~— не~функция среднего, а~функционал распределения, причём у~самого порога, где распределение бимодально из-за бистабильности редокс-переключателя. Усреднять там нечего.

\textbf{Динамическое среднее поле (DMFT).} Формально \emph{точно}: связь всех со~всеми через скаляр есть бесконечная размерность, и~в~пределе $J\to\infty$
\[D(t)=\mathrm{P}\big(v:\;x(t;v)<x_{\text{пор}}\big)\]
без всякого приближения. Предел существует и~измерен: поправки реализованного разброса сетки сходятся как~$1/J$, и~экстраполяция по~двум последним точкам даёт 0,1101. Но~экономии DMFT сама не~даёт: она сводит задачу к~одноклеточной плюс квадратура по~$v$, и~вся стоимость переезжает в~число узлов квадратуры.

\textbf{Уровневое множество.} Дало~бы полное решение~— одну границу вместо $J$~клеток. Мертво при~живой защите, \emph{точно верно} при~замороженной (раздел~\ref{sec:band}).

\textbf{Борн~— Оппенгеймер.} Разделение масштабов настоящее: 63~часа против быстрой релаксации. На~быстром масштабе доля повреждённых замкнута, медленная переменная двигает границу~— но~двигает её \emph{неодинаково по~клеткам}, и~отсюда полоса.

\textbf{Многоконфигурационное приближение.} Нужный путь, и~число конфигураций теперь известно из~измерения: три, плюс две подвижные границы.

\begin{note}
\textbf{Решение по~применению: редукцию не~делать ускорением.} Выигрыш оценивается втрое-вчетверо (32~узла против 8–12) против прогона, который уже укладывается в~три часа. Цена~— правка модели, не~проверяемая тождественностью, обнуляющая базу вердиктов и~требующая своего исследования сходимости по~$J$. И~главное: редукция огрубила~бы именно пограничный слой, то~есть единственное новое явление. Что окупается сразу и~бесплатно~— замкнутая уровневая формула при~замороженной защите как источник \emph{аналитических ожиданий}: модель не~трогается, кэш не~обесценивается, а~пороги получают вывод из~матожиданий и~дисперсий.
\end{note}

% =====================================================================
\section{Рваность нагрузки: гипотеза, которая не~прошла}\label{sec:burst}
Два расхождения с~арифметикой опасности смотрели в~\emph{противоположные} стороны, и~это выглядело зацепкой. Наблюдаемый прирост частоты срывов от~всплесков оказался вдесятеро \emph{сильнее} предсказанного; наблюдаемая доля повреждавшихся при~двойном суточном размахе~— втрое \emph{слабее}. Ни~выпуклость опасности, ни~замена плотности на~вероятность превышения такого дать не~могут: обе поправки двигают предсказание в~одну сторону. Оставался механизм, которого в~опасности нет вовсе,~— запаздывание защиты: короткий всплеск приходит на~неадаптированную клетку, а~поднятый на~сутки размах защита успевает частично скомпенсировать.

\textbf{Как это проверялось без единого подогнанного числа.} Дописать в~предсказатель множитель адаптации значило~бы ввести параметр, выводимый только из~алгебры складки, то~есть на~деле подогнать его. Поэтому вопрос был поставлен сравнением, в~котором всё лишнее закреплено: сила медленного шума задаёт \emph{дисперсию} нагрузки при~неизменном среднем (среднее логнормали скомпенсировано в~самой модели), а~время корреляции~— \emph{временной масштаб} при~неизменных среднем и~дисперсии. Перебор по~времени корреляции изолирует ровно одну вещь: успевает ли защита адаптироваться.

\textbf{Правило как записано~— не~выполнено.} Доля повреждавшихся по~упорядоченным временам корреляции дала $z=0{,}61$, $p=0{,}27$: тренда нет.

\textbf{И~дефект здесь мой, а~не~модельный.} Доля повреждавшихся~— «был~ли набор повреждён хоть раз за~60~суток»~— \emph{насыщается}: она идёт $0{,}599\to0{,}623\to0{,}626\to0{,}627\to0{,}614\to0{,}588$, то~есть горбом, и~чувствительности почти не~имеет. А~частота срывов на~тех~же блоках падает монотонно с~0,0550 до~0,0328~— на~40~\%,— и~длительность эпизода растёт с~5,37 до~8,04~суток (рис.~\ref{fig:v5}б). Прикидка по~опубликованной таблице даёт и~количественное согласие: доля падения, приходящаяся на~участок между двумя постоянными защиты, равна 0,523 против 0,424, ожидаемых при~равномерном по~логарифму падении.

\begin{note}
Эта прикидка \textbf{не является вердиктом} и~в~журнал не~внесена. Она посчитана после того, как данные были увидены, значит порог к~ним подогнан. Чтобы Ф24 получила честный вердикт, правило переписывается на~частоту срывов и~перезапускается на~\emph{другой} конфигурации~— другое зерно и~другая сила медленного шума,~— с~порогами, записанными заранее. Вердикт «не~выполнено» от~17.10 остаётся в~журнале как есть; дефект отмечен рядом. Это та~же дисциплина, по~которой мы не~переписываем задним числом ничего, включая собственные промахи.
\end{note}

\textbf{Урок шире самого правила.} К~четырём статистическим ловушкам версии~4 добавилась пятая, и~она о~выборе \emph{измеряемой величины}, а~не~критерия: насыщающаяся величина гасит настоящий эффект, и~никакая статистика этого не~исправит. Механизм действует на~частоту событий~— мерить надо частоту, а~не~долю тех, с~кем событие случилось хоть раз.

% =====================================================================
\section{Как мы теперь считаем}\label{sec:method}
Правило фиксируется до~счёта; это было с~самого начала. Версия~4 добавила к~нему вторую половину: \emph{у~каждого критерия должны быть известны его нулевое значение, его шум при~нашем числе наборов и~его мощность}. Иначе вердикт говорит о~выборке, а~не~о~модели.

Формально так. Часть правила считается проверяемой, только если расстояние от~нулевого значения до~порога не~меньше минимального различимого эффекта
\[\mathrm{MDE}=(z_{0{,}95}+z_{0{,}80})\,\mathrm{SE}=2{,}4865\,\mathrm{SE}.\]
Иначе вердикт~— «непроверяемо», и~\textbf{симметрично}: и~там, где правило провалилось~бы, и~там, где прошло~бы. Рядом печатается, сколько наборов нужно, чтобы порог стал различим. Для~утверждений вида «наблюдение согласуется с~предсказанием» допуск взят в~3~SE: он обязан быть шире MDE, иначе такое утверждение неопровержимо в~принципе.

Это не~формальность: четыре правила оказались непроверяемыми при~любом числе наборов, потому что мерились по~восьми ячейкам, где ошибка ранговой связи равна $1/\sqrt{7}=0{,}38$, а~порог ставился $-0{,}5$. Лечится это не~выборкой, а~дроблением ячеек.

\textbf{Четыре статистические ловушки, найденные по~дороге}, каждая стоила отдельного прогона:
\begin{enumerate}
\item \emph{Медиана по~наборам от~доли редких событий} всегда 0 или~1: она обнуляется, едва события есть меньше чем у~половины наборов. Три правила падали на~этом, а~не~на~модели.
\item \emph{Доля коротких промежутков} тождественно равна $1+\mathrm{CV}^2$ по~частотам наборов и~физики не~видит вовсе. Замена: $\mathrm{CV}^2$ промежутков \emph{внутри} набора, где при~глухом периоде $\mathrm{CV}^2=(1-t_r/g)^2$, а~нуль равен единице ровно.
\item \emph{Направление знакового критерия}: одна и~та~же функция, применённая к~«меньше~— лучше» и~«больше~— лучше», дала $p=0{,}99$ там, где правильное направление давало 0,038. Теперь направление~— обязательный аргумент, и~самопроверка гоняет оба.
\item \emph{Нулевая ошибка по~построению}: столбец «в~долях собственного здорового уровня» в~здоровой ячейке равен единице у~каждого набора, разброс нулевой, и~любой вердикт объявлялся достоверным.
\end{enumerate}

\textbf{Воспроизводимость.} Счёт разнесён по~файлам: модель, статистики, правила, сборщик блоков. Каждое правило \emph{объявляет}, какие блоки ему нужны; одинаковые блоки разных правил склеиваются по~хешу и~считаются один раз (из~78 запрошенных считается 42). Результат блока кладётся в~кэш, в~ключ которого входит хеш исходного кода модели: правка модели сама обесценивает кэш, и~старые блоки не~подмешаются к~новым. Отдельный тест требует \emph{побитового} совпадения с~прежним монолитом на~закреплённых точках. Вердикты пишутся в~журнал с~датой и~хешами блоков и~после счёта не~пересматриваются: негодность правила отмечается рядом, а~не~исправляется задним числом.

\begin{ineq}
Полезная мелочь, которая здесь дорого стоила. Размах $n$~значений величины с~относительной ошибкой~$r$ сам по~себе есть случайная величина: для~трёх независимых значений $\mathrm{E}[\text{размах}]=1{,}693\,r$. Поэтому порог «разброс по~$J$ не~больше 10~\%» осмыслен только при~$r<0{,}059$. При~300~наборах $r=12$~\%, то~есть правило было непроверяемо; при~2000 оно стало проверяемым~— и~тут~же не~прошло, обнаружив третий путь, которым $J$ входит в~модель.
\end{ineq}

% =====================================================================
\section{Точность счёта и~проверяемость утверждений}\label{sec:prec}
Версия~5 втрое ускорила счёт и~заодно выяснила, чему в~нём можно верить.

\textbf{Откуда взялось ускорение.} Счёт упирался не~в~арифметику, а~в~полосу памяти: каждая элементарная операция над~массивом клеток~— отдельный проход, и~шаг делал таких проходов около сорока пяти. Весь поклеточный блок переписан одним слитым циклом, где промежуточные величины живут в~регистрах. Выигрыш на~рабочем размере~— \textbf{3,3~раза по~стене}: полный прогон 123~блоков при~2400~наборах занимает 3~часа 05~минут вместо прежних почти четырёх часов при~меньшем объёме. Механизм виден на~рис.~\ref{fig:v5}г: у~счёта по~одной операции стоимость строки удваивается между 2000 и~5250~строк~— рабочий набор перестаёт влезать в~кэш,— и~дальше держится, а~у~слитого ядра она плоская на~всех размерах, потому что его рабочий набор~— одна строка.

\textbf{Чем доказано, что это та~же модель.} Быстрый путь не~переписан руками, а~\emph{собран из~исходника} эталона подстановкой текста с~обязательной проверкой единственности каждого куска: при~любой правке модели сборка падает, а~не~расходится молча. Сверка идёт на~двух уровнях~— все 63~поля результата на~пяти точках (включая обе стороны асимметрии; счётные совпадают побитово, худшее расхождение непрерывных $9{,}2\cdot10^{-15}$) и~отдельно весь набор правил, прогнанный дважды с~диффом всех вердиктов.

\textbf{Чему верить в~шаге по~времени.} Величины разделились на~два класса. Непрерывные средние сошлись с~запасом: при~вдвое более грубом шаге они держатся в~пределах 0,5~\%. Счётные~— нет: частота срывов падает монотонно во~всех условиях при~огрублении шага, и~экстраполяция Ричардсона даёт предел \emph{на~2~\% выше} значения при~рабочем шаге. Причина не~в~моменте пересечения порога: подшаговая интерполяция его уточняет, но~даёт всего 0,009~\%. Смещение берётся из~пропущенных подшаговых экскурсий и~из~погрешности самой траектории. Вывод для~чтения результатов: порогам по~счётным величинам тоньше примерно 3~\% верить нельзя, и~это записано как известный предел, а~не~замазано.

\textbf{Седьмая ловушка, и~она однажды уже попала в~текст этой брошюры.} \emph{Среднее по~бимодальной популяции.} Наборы в~рабочей точке распадаются на~две кучи: либо сломаны необратимо (доля повреждённых выше 0,9), либо не~тронуты (ниже 0,01), а~между ними почти никого. Среднее по~такой популяции~— это \emph{доля сломанных наборов}, а~не~состояние системы, и~читать его как состояние нельзя. Именно так в~версии~4 возникло утверждение, будто модель не~достаёт до~лестницы захвата: сравнивалось среднее по~смеси, тогда как у~сломанных наборов повреждение и~тогда стояло у~единицы (раздел~\ref{sec:bifur}). Правило чтения: всюду, где в~таблицах встречается доля повреждённых, рядом надо держать долю наборов выше 0,9 и~ниже 0,01, иначе среднее обманет.

\textbf{Пятая и~шестая ловушки.} Пятая~— выбор измеряемой величины (раздел~\ref{sec:burst}). Шестая~— \emph{неопровержимое утверждение о~согласии}: допуск равносильности, превышающий собственный размах кривой, пропустит любые данные. Теперь правило обязано проверять и~это, и~при~нарушении выдаёт «непроверяемо» с~указанием, сколько наборов нужно.

\begin{ineq}
Отдельный приём, окупившийся сразу. Пять частей правил были непроверяемы не~по~нехватке наборов, а~по~\emph{информационному пределу}: ранговая связь мерилась по~восьми ячейкам, её ошибка $1/\sqrt{n-1}=0{,}38$, то~есть различимый порог 0,94 против записанного~$-0{,}50$. Мельчить ячейки нечего: кривая индексируется бодрствованиями, их около одного в~сутки. Выход~— считать наклон у~\emph{каждого набора} по~его ячейкам и~проверять разброс между наборами, где $n$~— сотни. Порог при~этом обязан быть пересчитан: кривая одного набора шумнее сводной, и~связь ослаблена классическим множителем
\[\rho_{\text{набора}}=\rho_{\text{сводной}}\big/\sqrt{1+\lambda},\qquad \lambda=\frac{\text{дисперсия шума ячейки}}{\text{дисперсия сигнала}},\]
причём $\lambda$ оценивается из~тех~же данных \emph{до} того, как мы посмотрели на~наклоны. На~синтетике с~известным ответом формула попадает с~точностью 1–5~\%. В~деле $\lambda=33{,}8$, порог $-0{,}50$ пересчитался в~$-0{,}085$, а~MDE упал с~0,94 до~0,076~— в~двенадцать раз,— и~четыре части из~пяти закрылись на~уже посчитанных данных.
\end{ineq}

\textbf{Два правила перестроены по~ходу.} Правило о~составе эпизода жило на~\emph{втором} интеграторе, копии шага внутри самого правила, со~своей самосверкой против эталона. Сверка отказала, едва в~эталон попала интерполяция: расхождение вышло ровно в~один шаг на~каждую границу эпизода. Самосверка сработала как задумано, но~вывод из~этого один~— дубликат расходится при~каждой правке модели. Он удалён, записи эпизодов пишет сам эталон, и~расходиться больше нечему. Попутно выборка выросла в~восемь раз: 438~наборов и~5323~эпизода вместо 62 и~788. Второе~— разметка бодрствований: они попадали в~ячейку по~\emph{максимуму} повреждения за~это~же бодрствование, то~есть метка частично создавалась измеряемой длиной. Обратная причинность снята разметкой по~состоянию на~\emph{начале} бодрствования.

% =====================================================================
\section{Что мы считаем бесспорным}\label{sec:solid}
К~версии~5 журнал вердиктов содержит 106~частей правил: 57~выполнено, 39~не~выполнено, 9~непроверяемо, одна ниже вилки. Ниже~— та часть, которую мы готовы защищать без~оговорок: утверждения с~порогом, назначенным до~счёта из~матожиданий и~дисперсий, с~парным контролем и~с~проверкой на~опровержимость.

\textbf{Про механизм.}
\begin{enumerate}[leftmargin=1.3em,itemsep=2pt,topsep=2pt]
\item Длительность эпизода~— \emph{прямая} функция памяти защиты, без~насыщения; проверочная точка вне подгонки подтвердила форму. Отсюда: пассивное усиление ремонта клинической цели не~достигает, промах в~2–10~раз.
\item Разброс длительности эпизодов внутри набора на~99,8~\% объясняется ходом нагрузки и~на~0,2~\%~— починкой. Разложение точное по~построению, сумма долей сходится до~$10^{-16}$.
\item Длинная память защиты делает предвестник: при~её отнятии высота падает у~107~наборов из~174, $p=0{,}0015$. Медленный шум нагрузки для~этого не~нужен~— его обоснование снято дважды.
\item Асимметрия постоянных подъёма и~спада~— режимный параметр, а~не~деталь: средний уровень защиты следует выкладке $w(r)$ с~точным совпадением порядка, $p=0{,}0083$ (раздел~\ref{sec:asym}).
\item Популяция живёт в~трёх фазах. Полоса спасённых существует (семнадцать сигм), обращается в~\emph{ровный нуль} при~замороженной защите и~расширяется вдвое при~прайминге (раздел~\ref{sec:band}).
\item Ритм ломается лестницей захвата с~рациональными числами вращения и~с~чередованием длин циклов на~дробных ступенях; каскад идёт прибавлением единицы.
\end{enumerate}

\textbf{Про метод.}
\begin{enumerate}[leftmargin=1.3em,itemsep=2pt,topsep=2pt,resume]
\item Число клеток входит в~модель как эффективный разброс стойкости, а~не~как разрешение счёта; предел $J\to\infty$ существует и~измерен.
\item Шесть статистических ловушек, каждая из~которых однажды дала~бы ложный вывод, и~каждая закрыта проверкой в~коде.
\item Тождественность ускоренного счёта доказана на~двух уровнях~— побитово по~полям и~по~всем вердиктам сразу.
\end{enumerate}

\textbf{Чего мы \emph{не} утверждаем.} Форма зависимости ширины полосы от~доли повреждённых~— опровергнута дважды с~противоположными знаками. Объяснение расхождений арифметики опасности запаздыванием защиты~— не~подтверждено в~записанной форме; переформулировка на~частоту срывов выглядит состоятельной, но~вердикта пока не~имеет и~ждёт прогона на~независимой конфигурации. Поломка ритма при~\emph{откалиброванной} связи «повреждение~— сон» недостижима, и~то, что модель доходит до~лестницы при~вдвое более сильной связи, само по~себе ничего о~болезни не~говорит. Количественная форма $w(r)$, девять частей со~статусом «непроверяемо» и~все величины, чьи пороги тоньше 3~\% по~счётным мерам,~— вне круга бесспорного.

% =====================================================================
\section{Открытые вопросы}\label{sec:crit}
\begin{enumerate}
\item \textbf{Одна клетка вместо популяции.} Повреждение одной клетки перекладывает нагрузку на~соседние; это может и~распространять повреждение, и~сдерживать его.
\item \textbf{Параметры не~от~тех клеток.} Кинетика пероксиредоксина измерена в~клеточных линиях. Вывод «клетки стоят у~границы окна» относится к~ним, а~не~к~PV-клеткам.
\item \textbf{Что такое нагрузка.} Часть притока окислителей постоянна (митохондрии), часть зависит от~работы рецепторов и~приходит от~соседних клеток. Доли неизвестны.
\item \textbf{Нет памяти между воздействиями.} \emph{Закрыто в~версии~4:} память появилась дважды~— как двухступенчатая защита с~постоянной 63~ч и~как прайминг (разделы~\ref{sec:slow} и~\ref{sec:prime}). Опыт с~двумя дозами кетамина теперь есть чем описывать.
\item \textbf{Суточный ритм: навязан или собственный.} Либо часы периодически меняют нагрузку, либо переключатель колеблется сам и~лишь подстраивается. Во~втором случае болезнь — это остановка цикла в~повреждённой фазе.
\item \textbf{Сдвиг фаз больше четверти периода} требует ещё одной медленной стадии или звена дробного порядка.
\item \textbf{Спектр медленных колебаний нагрузки} на~минутах и~часах не~измерен, а~от~него зависит, может~ли шум перебросить систему.
\item \textbf{Гибель или молчание.} Данные о~больных говорят скорее о~молчании клеток.
\item \textbf{Нет развития, дофамина и~симптомов.} Модель описывает состояние клетки, а~не~болезнь.
\item \textbf{Какое звено первично.} Окислительная защита, сигнал поддержки BDNF или оба — следствия общей причины. От~ответа зависит, что считать медленной переменной и~какие времена подставлять.
\item \textbf{Третья медленная величина.} \emph{Наполовину закрыто в~версии~5:} несимметрия подъёма и~спада защиты теперь в~модели, измерена и~ведёт себя по~выкладке $w(r)$ (раздел~\ref{sec:asym}), так что этот кандидат больше не~гипотеза. Остаётся второй: предвестник сна объясняется памятью защиты только на~пятую часть, и~времён, кроме уже заложенных, в~уравнениях нет. Либо это заполнение популяции (у~него нет своего~$\tau$), либо нужна ещё одна медленная стадия.
\item \textbf{Что задаёт свободный член закона длительности.} Прямая $1{,}28+0{,}97\,\tau_{\rm б}$ при~нулевой памяти даёт 1,3~суток, и~объяснить этот остаток возвратом~$x$ нельзя: он возвращается за~часы.
\item \textbf{Поломка ритма недостижима при~рабочей связи.} Лестница захвата существует, и~модель до~неё \emph{доходит}~— но~при~связи «повреждение~— сон» не~ниже 4,77~мВ, то~есть в~2,4~раза сильнее рабочих 2~мВ. Второй выход, записанный в~версии~4 (хроническое повреждение, держащее долю сломанных высокой), \emph{закрыт измерением}: тождественно $kD\le k$, а~у~сломанных наборов доля уже на~потолке. Остаётся одно из~двух: либо связь в~модели слишком слаба, либо поломка ритма при~болезни делается не~этим путём.
\item \textbf{Сама идея двух состояний.} Прямого опыта, показывающего гистерезис в~PV-клетке, нет; опыт с~двумя дозами кетамина — ближайшее к~нему наблюдение.
\item \textbf{Чем задана ширина полосы спасённых.} Геометрия пограничного слоя опровергнута: два перебора дают противоположные знаки (раздел~\ref{sec:band}). Значит ширину определяет не~наблюдаемая доля повреждённых, а~что-то третье~— вероятнее всего время, проведённое популяцией у~самого опрокидывания. Проверяемо против уже измеряемой доли времени в~середине шкалы, и~обязано совпасть на~\emph{обоих} переборах, иначе снимается так~же.
\item \textbf{Почему короткое воздействие бьёт сильнее длинного.} Запаздывание защиты в~записанной форме не~подтвердилось (раздел~\ref{sec:burst}); переформулировка на~частоту срывов выглядит состоятельной, но~вердикта не~имеет. Если и~она не~пройдёт, следующий кандидат~— нелинейность за~верхом окна нагрузки: всплеск выходит за~него, а~суточный размах нет, и~закон опасности там может быть другим.
\item \textbf{Сколько стоит правда о~счётных величинах.} Частота срывов при~рабочем шаге занижена примерно на~2~\%, и~это предел выборки, а~не~округления (раздел~\ref{sec:prec}). Либо шаг вдвое мельче и~вдвое дороже, либо точное восстановление подшаговых экскурсий из~долей пересечения, которые в~слитом ядре уже лежат в~регистрах.
\end{enumerate}

% =====================================================================
\section{Словарь}
{\small
\begin{description}[leftmargin=1.2em,itemsep=1pt,parsep=0pt,font=\normalfont\bfseries]
\item[Аксон] выходной отросток нейрона
\item[AMPA-рецептор] быстрый рецептор глутамата
\item[BDNF] фактор роста нейронов; сигнал поддержки, который выделяют активные клетки
\item[Астроцит] опорная клетка мозга; снабжает нейроны, в~том числе сырьём для глутатиона
\item[АТФ] молекула, в~которой клетка запасает энергию
\item[Белок] цепочка аминокислот, собранная по~гену; основной рабочий материал клетки
\item[ГАМК] главный тормозной медиатор
\item[Гамма-ритм] колебания активности коры с~частотой 30–80~Гц
\item[GAD67] фермент, производящий ГАМК
\item[Гиппокамп] отдел мозга, связанный с~памятью; в~схеме Грейса — источник избыточного сигнала
\item[Главные часы] отдел мозга над перекрёстом зрительных нервов (супрахиазматическое ядро), синхронизирующий суточные ритмы тела
\item[Глутамат] главный возбуждающий медиатор
\item[Глутатион] главный расходуемый антиоксидант клетки
\item[Дендрит] входной отросток нейрона
\item[Дофамин] медиатор, связанный с~оценкой значимости событий
\item[Интернейрон] тормозной нейрон коры
\item[Кататония] состояние с~резкими двигательными нарушениями; в~МКБ-11 — отдельный диагноз
\item[Кетамин] вещество, блокирующее NMDA-рецептор
\item[Культура клеток] клетки, выращенные вне организма
\item[Медиатор] молекула, передающая сигнал через синапс
\item[Мембрана] оболочка клетки
\item[МКБ-11] международная классификация болезней последнего пересмотра
\item[Митохондрия] «электростанция» клетки; побочно производит окислители
\item[мРНК] рабочая копия гена, по~которой строится белок
\item[NMDA-рецептор] медленный рецептор глутамата, пропускает кальций; датчик активности
\item[Nox2] фермент, намеренно производящий окислители
\item[Окислитель] молекула, отнимающая электроны у~других молекул
\item[Окислительный стресс] избыток окислителей над возможностями очистки
\item[Парвальбумин (PV)] белок, связывающий кальций; метка быстрых тормозных клеток
\item[Перинейрональная сеть] оболочка из~внеклеточного вещества вокруг PV-клетки
\item[Пероксиредоксин] главный фермент-уловитель перекиси
\item[Пирамидная клетка] возбуждающий нейрон коры
\item[Продром] период перед первым эпизодом психоза
\item[Психоз] состояние с~галлюцинациями и~бредом
\item[Растормаживание] рост активности из-за ослабления торможения
\item[Рецептор] белок, принимающий сигнал; у~нейронов — ворота для ионов, открываемые медиатором
\item[Серотонин] медиатор, управляющий общим режимом работы коры; на~его рецепторы действуют галлюциногены
\item[Синапс] место контакта двух нейронов
\item[Сульфиредоксин] фермент, чинящий переокисленный пероксиредоксин
\item[TrkB] рецептор, принимающий сигнал BDNF
\item[Фермент] белок, ускоряющий химическую реакцию
\item[Экспрессия] производство белка по~гену
\end{description}}

% =====================================================================
\section{Источники}
{\small
\textbf{Сверено при подготовке брошюры} (метка \tP):
\begin{itemize}
\item Saha~S. и~соавт. A systematic review of the prevalence of schizophrenia. PLoS Medicine, 2005.
\item Hilker~R. и~соавт. Heritability of schizophrenia and schizophrenia spectrum based on the nationwide Danish twin register. Biological Psychiatry, 2018 (по~изложению пресс-релиза).
\item Krystal~J.\,H. и~соавт. Subanesthetic effects of the noncompetitive NMDA antagonist, ketamine, in humans. Archives of General Psychiatry, 1994.
\item Kapur~S. и~соавт. Relationship between dopamine D2 occupancy, clinical response, and side effects. American Journal of Psychiatry, 2000.
\item Howes~O.\,D. и~соавт. The nature of dopamine dysfunction in schizophrenia and what this means for treatment. Archives of General Psychiatry, 2012.
\item Работа об~обороте глутатиона в~мозге мыши, запись PubMed 9165485.
\item Harkness~J.\,H. и~соавт. Diurnal changes in perineuronal nets and parvalbumin neurons in the rat medial prefrontal cortex. Препринт bioRxiv, 2020; опубликовано в~2021.
\item Molecular and cellular rhythms in excitatory and inhibitory neurons in the mouse prefrontal cortex. Препринт bioRxiv, 2024.
\item Schizophrenia and catatonia: from ICD-10 to ICD-11. Обзорная статья, Springer, 2025 (авторы не~сверены). Перечень кодов группы 6A2 МКБ-11 — по~справочнику кодов.
\item Hashimoto~T. и~соавт. Relationship of brain-derived neurotrophic factor and its receptor TrkB to altered inhibitory prefrontal circuitry in schizophrenia. Journal of Neuroscience, 2005.
\item Fernandes~B.\,S. и~соавт. Peripheral brain-derived neurotrophic factor in schizophrenia and the role of antipsychotics: meta-analysis and implications. Molecular Psychiatry, 2015.
\item Durany~N., Thome~J. Neurotrophic factors and the pathophysiology of schizophrenic psychoses. European Psychiatry, 2004; Angelucci~F. и~соавт. BDNF in schizophrenia, depression and corresponding animal models. Molecular Psychiatry, 2005.
\item Jaggar~M., Vaidya~V.\,A. 5-HT2A receptors and BDNF regulation: implications for psychopathology. Глава в~книге, Springer, 2018 (по~аннотации).
\item Papadia~S. и~соавт. Synaptic NMDA receptor activity boosts intrinsic antioxidant defenses. Nature Neuroscience, 2008 (по~пересказу).
\end{itemize}
\textbf{Прочитано авторами проекта} (метка \tS):
\begin{itemize}
\item Hardingham~G.\,E., Do~K.\,Q. Linking early-life NMDAR hypofunction and oxidative stress in schizophrenia pathogenesis. Nature Reviews Neuroscience, 2016.
\item Cuenod~M. и~соавт. Обзор о~порочных кругах через окислительный стресс, 2021, и~ответ на~критику (Molecular Psychiatry, 2022).
\item Selvaggio~G., Coelho~P., Salvador~A. Модель системы пероксиредоксин–тиоредоксин. Redox Biology, 2018.
\item Volman~V., Behrens~M.\,M., Sejnowski~T.\,J. Модель сети с~потерей парвальбумина и~GAD67, 2011.
\item Kinney~J.\,W. и~соавт., 2006 — кетамин в~культуре; Behrens~M.\,M. и~соавт., 2008 — кетамин у~мышей; Baxter~P.\,S. и~соавт., 2015 — рецептор и~синтез глутатиона; Brennan~A.\,M. и~соавт., 2009 — рецептор и~Nox2.
\item Kil~I.\,S. и~соавт. Circadian oscillation of sulfiredoxin in the mitochondria. Molecular Cell, 2015; del~Olmo~M. и~соавт. A robust model for circadian redox oscillations. International Journal of Molecular Sciences, 2019.
\item Edgar~R.\,S. и~соавт., 2012 — суточные циклы пероксиредоксинов; Wang~T.\,A. и~соавт., 2012 — окислительное состояние и~возбудимость в~главных часах; Kinoshita~C. и~соавт., 2014 — суточный ритм глутатиона.
\item Harkness~J.\,H. и~соавт. — нарушение сна и~окислительный стресс в~PV-клетках; Vyazovskiy~V.\,V. и~соавт., 2009 — частота разрядов и~сон; He~B.\,J., 2014 — спектр активности мозга без выделенного масштаба.
\item Работы группы A.~Grace о~модели MAM; работы O.~Howes о~синтезе дофамина перед психозом; Do~K.\,Q. и~соавт., 2000, и~сводка 2021~года о~глутатионе у~больных.
\item Обзор о~строении и~работе PV-клеток; работы об~устойчивости коры (Tsodyks, 1997; Ozeki, 2009); исследования длительности симптомов и~продрома.
\end{itemize}
Выходные данные статей второй группы приведены по~рабочим записям и~требуют сверки перед цитированием.

\textbf{Учебные знания} (метка \tU) записаны по~памяти. Перед тем как ссылаться на~такой факт, его следует сверить по~учебнику, например: Kandel~E. и~соавт., Principles of Neural Science; Purves~D. и~соавт., Neuroscience; Alberts~B. и~соавт., Molecular Biology of the Cell.
}
\end{document}
